% =====================================================================
% Gap mechanisms, Section 2: the local resonator gap
% Transcription of Dylan's handwritten notes (dispersion_theory_page1.pdf,
% dispersion_theory_page2.pdf, dispersion_theory_pages3-4.pdf), Nov 2026.
% Corrections of the handwritten text are set in blue (\fixed{...}) and
% explained in a red margin note (\corr{...}).
% Requires: amsmath, amssymb, graphicx, xcolor. The macros \fixed and \corr are
% defined below with \providecommand, so this file works in any main file that
% loads xcolor (a main file may define its own versions first).
% Include in the main file with  % =====================================================================
% Gap mechanisms, Section 2: the local resonator gap
% Transcription of Dylan's handwritten notes (dispersion_theory_page1.pdf,
% dispersion_theory_page2.pdf, dispersion_theory_pages3-4.pdf), Nov 2026.
% Corrections of the handwritten text are set in blue (\fixed{...}) and
% explained in a red margin note (\corr{...}).
% Requires: amsmath, amssymb, graphicx, xcolor. The macros \fixed and \corr are
% defined below with \providecommand, so this file works in any main file that
% loads xcolor (a main file may define its own versions first).
% Include in the main file with  % =====================================================================
% Gap mechanisms, Section 2: the local resonator gap
% Transcription of Dylan's handwritten notes (dispersion_theory_page1.pdf,
% dispersion_theory_page2.pdf, dispersion_theory_pages3-4.pdf), Nov 2026.
% Corrections of the handwritten text are set in blue (\fixed{...}) and
% explained in a red margin note (\corr{...}).
% Requires: amsmath, amssymb, graphicx, xcolor. The macros \fixed and \corr are
% defined below with \providecommand, so this file works in any main file that
% loads xcolor (a main file may define its own versions first).
% Include in the main file with  % =====================================================================
% Gap mechanisms, Section 2: the local resonator gap
% Transcription of Dylan's handwritten notes (dispersion_theory_page1.pdf,
% dispersion_theory_page2.pdf, dispersion_theory_pages3-4.pdf), Nov 2026.
% Corrections of the handwritten text are set in blue (\fixed{...}) and
% explained in a red margin note (\corr{...}).
% Requires: amsmath, amssymb, graphicx, xcolor. The macros \fixed and \corr are
% defined below with \providecommand, so this file works in any main file that
% loads xcolor (a main file may define its own versions first).
% Include in the main file with  \include{gap_mechanisms_local_resonance}
% =====================================================================

\providecommand{\fixed}[1]{{\color{blue!70!black}#1}}
\providecommand{\corr}[1]{\marginpar{\raggedright\footnotesize\color{red!70!black}#1}}

\chapter{Local resonator gap}
\label{ch:local_resonator_gap}

Let $m$ be the point mass, connected at $\mathbf{x}_0$ to the plate by a spring
of stiffness $k_s$. Write the plate deflection at the attachment point as
$w_0 := w(\mathbf{x}_0)$ and the displacement of the resonator mass as $u$.
The spring is stretched by $w_0 - u \Rightarrow m\ddot{u} = k_s(w_0 - u)$.
Assuming a harmonic motion $\ddot{u} = -\omega^2 u$, we obtain
\begin{equation*}
  -m\omega^2 u = k_s(w_0 - u)
  \;\Rightarrow\;
  u = \frac{k_s w_0}{k_s - m\omega^2}
  \;\overset{\cdot\frac{1/m}{1/m},\ \omega_r^2 = \frac{k_s}{m}}{=}\;
  \frac{\omega_r^2 w_0}{\omega_r^2 - \omega^2} .
\end{equation*}
The spring pulls on the plate with $F = k_s(u - w_0)$ according to Newton~III.
We find
\begin{equation*}
  F = k_s(u - w_0)
  \overset{k_s = m\omega_r^2}{=}
  m\omega_r^2\left(\frac{\omega_r^2}{\omega_r^2 - \omega^2} - 1\right) w_0
  = m\omega_r^2 \cdot \frac{\omega^2}{\omega_r^2 - \omega^2}\, w_0
  =: m_\mathrm{eff}(\omega)\,\omega^2 w_0 ,
\end{equation*}
where we have defined the \underline{effective mass}
$m_\mathrm{eff}(\omega) := m\,\dfrac{\omega_r^2}{\omega_r^2 - \omega^2}$.
Everything from Section~1, \emph{Bragg gaps}, now carries over by making the
replacements $m \to m_\mathrm{eff}(\omega)$ and
$\mu \to \mu_\mathrm{eff}(\omega) := \mu\,\dfrac{\omega_r^2}{\omega_r^2 - \omega^2}$.

We discuss different cases:
\begin{enumerate}
  \item $\omega \ll \omega_r$: $m_\mathrm{eff} \approx m$, $u \approx w_0$, like a
        rigid mass.
  \item $\omega$ just below $\omega_r$: $m_\mathrm{eff} \to \infty$,
        $u \gg w_0$, in phase, plate sees a huge mass and is pinned.
  \item $\omega$ just above $\omega_r$: $m_\mathrm{eff} \to -\infty$,
        $u \gg w_0$, antiphase, spring pulls plate back against its motion.
  \item $\omega \gg \omega_r$: $m_\mathrm{eff} \approx -m\,\dfrac{\omega_r^2}{\omega^2}$,
        $u \approx 0$, $F \to -k_s w_0 \mathrel{\widehat{=}}$ a spring to a fixed point.
\end{enumerate}

When considering \underline{waves much larger than the unit cell dimension $a$},
$m_\mathrm{eff}$ can be considered spread over the entire cell area:
$\rho h \to \rho h + m_\mathrm{eff}/A$.
Consider the Kirchhoff relation
\begin{equation*}
  D k^4 = \rho h\,(1 + \mu_\mathrm{eff})\,\omega^2
  = \rho h\left(1 + \mu\,\frac{\omega_r^2}{\omega_r^2 - \omega^2}\right)\omega^2
  = \rho h\,\frac{\omega_r^2(1+\mu) - \omega^2}{\omega_r^2 - \omega^2}\,\omega^2 .
\end{equation*}
A real solution needs a positive RHS, meaning that a negative RHS corresponds to
a band gap. As $\rho h, \omega^2 > 0$, consider the fraction: It is negative if
both the numerator $> 0$ and the denominator $< 0$ \fixed{(the opposite
combination, numerator $< 0$ and denominator $> 0$, would need
$\omega_r\sqrt{1+\mu} < \omega < \omega_r$, which is impossible)},%
\corr{Added: the second sign combination and why it is excluded.}
namely
\begin{equation*}
  \omega_r^2(1+\mu) - \omega^2 \overset{!}{>} 0
  \;\Rightarrow\; \omega < \omega_r\sqrt{1+\mu}
  \quad\text{and}\quad
  \omega_r \overset{!}{<} \omega ,
  \quad\text{thus}
\end{equation*}
\corr{Handwritten box: $\omega_r < \omega < \omega_r(1+\mu)$; the square root
was missing.}%
\begin{equation}
  \boxed{\;\omega_r < \omega < \fixed{\omega_r\sqrt{1+\mu}}\;}
  \label{eq:effective_medium_gap}
\end{equation}
This is limited by the fact that only $\mathbf{G} = 0$ is kept (i.e.,
$\tau = f_r/f_X \to 0$) and even at $\Gamma$ (small $\mathbf{k} \Rightarrow$
long waves) the upper edge is too high, as other waves add a local bending sum
$S$, lowering the upper edge by about $0.051\,\tau^2$\fixed{$\,f_X$ (fit for
$\mu = 0.25$)}.%
\corr{The $0.051\,\tau^2$ fit holds for $\mu = 0.25$ only and is in units
of $f_X$. Exact at $f_r = 1200\,\mathrm{Hz}$: $0.056\,f_X$
($1341.6 \to 1274.7\,\mathrm{Hz}$).}
Figure~\ref{fig:dispersion_with_il} shows how far the effective-medium band
\eqref{eq:effective_medium_gap} is off for the production design.

\begin{figure}[htbp]
  \centering
  \includegraphics[width=\textwidth]{Figures/dispersion_cell0.1m_1200hz_mass0.25_0to5000hz_with_il.pdf}
  \caption{Left: dispersion diagram (FE, Bloch) of the $0.1\,\mathrm{m}$ unit
    cell ($h = 5\,\mathrm{mm}$, steel) with one resonator,
    $f_r = 1200\,\mathrm{Hz}$, $\mu = 0.25$, along
    $\Gamma$--$X$--$M$--$\Gamma$; the color gives the plate share of the
    kinetic energy (blue: plate, red: resonator). The complete gap runs from
    $1031$ to $1195\,\mathrm{Hz}$ ($14.7\,\%$) and lies \emph{below} $f_r$,
    not in the effective-medium band
    $f_r \dots f_r\sqrt{1+\mu} = 1200\dots1342\,\mathrm{Hz}$ of
    Eq.~\eqref{eq:effective_medium_gap}: $f_r$ sits on the first Bragg fold
    ($f_X = 1194.4\,\mathrm{Hz}$). The two candidates for the upper edge of
    Eq.~\eqref{eq:upper_edge_candidates} are both visible: the flat resonator
    band at $\Gamma$ ($1273.5\,\mathrm{Hz}$, the effective-medium value lowered
    by the bending sum $S$) and the sin pattern at $X$ ($f_X$), which is the
    lower one here and sets the edge. The lower edge is band~1 at $M$. Right:
    insertion loss of the radiated power of the $1 \times 1\,\mathrm{m}$ plate
    with $10 \times 10$ resonators on the same $0.1\,\mathrm{m}$ lattice (same
    $f_r$ and $\mu$), for equal force and for equal input power. The IL peaks
    inside the gap and drops to its minimum just above it, at the flat
    resonator band.}
  \label{fig:dispersion_with_il}
\end{figure}

Keeping only one plate mode we obtain $w(\mathbf{x}) = q\,\varphi(\mathbf{x})$,
where $q$ is the modal amplitude. Let $u$ be the resonator's displacement and
$k_s$ its stiffness. Again assuming a harmonic response, i.e.,
$\ddot{u} = -\omega^2 u = -\lambda u$, we find
\begin{equation}
  k_s(u - \varphi q) = \lambda m u
  \label{eq:resonator_row}
\end{equation}
for the resonator. For the plate mode we find in analogy to the rigid mass
\begin{equation*}
  (\mathcal{L} - \lambda\rho h)\, w = f(\mathbf{x}) ,
  \qquad f(\mathbf{x}) = F\,\delta(\mathbf{x} - \mathbf{x}_0) .
\end{equation*}
Inserting $w = q\varphi$, multiplying by $\varphi^*$ and integrating over the
cell we again obtain (omission of the argument $\mathrel{\widehat{=}}$ evaluated
at $\mathbf{x}_0$)
\begin{equation}
  \lambda_p q - \lambda q = k_s\varphi^* u - k_s|\varphi|^2 q
  \;\Rightarrow\;
  (\lambda_p + k_s|\varphi|^2)\, q - k_s\varphi^* u = \lambda q .
  \label{eq:plate_row}
\end{equation}
Writing \eqref{eq:resonator_row} and \eqref{eq:plate_row} into matrix form
$K\mathbf{x} = \lambda M\mathbf{x}$, $\mathbf{x} := (q, u)^T$ we find
\begin{equation*}
  K = \begin{pmatrix}
        \lambda_p + k_s|\varphi|^2 & -k_s\varphi^* \\
        -k_s\varphi & k_s
      \end{pmatrix} ,
  \qquad
  M = \operatorname{diag}(1, m) .
\end{equation*}
Transforming the resonator coordinate into a mass-normalized coordinate, i.e.,
$v := \sqrt{m}\,u$ and identifying
\begin{gather*}
  k_s|\varphi|^2 = \frac{k_s}{\rho h A} = \frac{k_s}{m}\cdot\frac{m}{\rho h A}
  = \mu\lambda_r ; \qquad \frac{k_s}{m} = \lambda_r ; \\
  \frac{k_s\varphi}{\sqrt{m}} = \frac{k_s}{\sqrt{m\rho h A}}
  = \frac{k_s}{m}\sqrt{\frac{m}{\rho h A}} = \sqrt{\mu}\,\lambda_r
\end{gather*}
we see that $K$ becomes a Hermitian matrix $H$
\begin{equation*}
  H = \begin{pmatrix}
        \lambda_p + \mu\lambda_r & -\sqrt{\mu}\,\lambda_r \\
        -\sqrt{\mu}\,\lambda_r & \lambda_r
      \end{pmatrix}
\end{equation*}
and $M$ becomes the identity. The eigenvalues at equal diagonals, i.e.,
$\lambda_p = (1-\mu)\lambda_r$, are
\begin{equation*}
  \lambda_{1,2} = (1 \pm \sqrt{\mu})\,\lambda_r .
\end{equation*}

Now we want to consider the influence of other modes. For this, we rewrite
\eqref{eq:resonator_row} using
$\dfrac{m\omega^2}{k_s} = \dfrac{\omega^2}{\omega_r^2} =: \beta$ to find
\begin{equation*}
  u = \frac{w_0}{1 - \beta} .
\end{equation*}
Newton~III then gives us for the force on the
\fixed{plate}\corr{Handwritten: ``force on the mass''. $F = m\omega^2 u$ is
the force the mass exerts on the plate; the force on the mass is
$-m\omega^2 u$.}
\begin{equation*}
  F = m\omega^2 u = m\omega^2\,\frac{w_0}{1-\beta} = m_\mathrm{eff}\,\omega^2 w_0
  \qquad (\mu_\mathrm{eff}\text{ analogous}) .
\end{equation*}

As before, we find in our known notation
\begin{equation*}
  (\lambda_n^0 - \lambda)\, q_n = F\varphi_n^*(\mathbf{x}_0)
  \;\Rightarrow\;
  q_n = F\,\frac{\varphi_n^*(\mathbf{x}_0)}{\lambda_n^0 - \lambda} ,
\end{equation*}
i.e., how strongly each wave responds to the point force. We then find
\begin{equation*}
  w_0 = \sum_n \varphi_n(\mathbf{x}_0)\, q_n
  = F \sum_n \frac{|\varphi_n(\mathbf{x}_0)|^2}{\lambda_n^0 - \lambda}
  =: C F ,
\end{equation*}
where we have defined the compliance $C$, the
\fixed{deflection per unit force}.%
\corr{Handwritten: ``the force per unit deflection''. That is $Z$;
$C = w_0/F$ is the deflection per unit force.}
Defining $Z := m_\mathrm{eff}\lambda$, the force per unit deflection, we finally
obtain
\begin{equation*}
  w_0 = w_0 Z C \;\Rightarrow\; \boxed{1 = ZC} \qquad (w_0 \neq 0) .
\end{equation*}
Using $|\varphi(\mathbf{x}_0)|^2 = \dfrac{1}{\rho h A}$ we get
\begin{equation}
  \boxed{\;1 = \mu_\mathrm{eff}\sum_n \frac{\lambda}{\lambda_n^0 - \lambda}\;}
  \label{eq:pwe_condition}
\end{equation}
the plane-wave expansion (PWE) condition. At every $\mathbf{k}$, every
$\lambda$ satisfying it is one dot in the diagram. For $w_0 = 0$ we find
\begin{equation*}
  (\lambda_n^0 - \lambda)\, q_n = 0 \;\Rightarrow\; \lambda_n^0 = \lambda ,
\end{equation*}
thus, these modes keep their bare frequencies.

With this, we may find the upper edge of the band gap. First, writing out the
PWE for one \fixed{single}%
\corr{Handwritten: ``singular''.}
plate frequency $\lambda_p$, using
\begin{equation*}
  1 - \beta = 1 - \frac{\lambda}{\lambda_r} = \frac{\lambda_r - \lambda}{\lambda_r}
\end{equation*}
we get
\begin{equation*}
  1 = \mu\,\frac{\lambda_r}{\lambda_r - \lambda}\cdot\frac{\lambda}{\lambda_p - \lambda}
  \;\Rightarrow\;
  (\lambda_p - \lambda)(\lambda_r - \lambda) = \mu\lambda_r\lambda
\end{equation*}
and computing the characteristic polynomial of $H$ we find
\begin{align*}
  (\lambda_p + \mu\lambda_r - \lambda)(\lambda_r - \lambda) - \mu\lambda_r^2
  &= \bigl[(\lambda_p - \lambda) + \mu\lambda_r\bigr](\lambda_r - \lambda) - \mu\lambda_r^2 \\
  &= (\lambda_p - \lambda)(\lambda_r - \lambda) + \mu\lambda_r(\lambda_r - \lambda) - \mu\lambda_r^2
   = (\lambda_p - \lambda)(\lambda_r - \lambda) - \mu\lambda_r\lambda \overset{!}{=} 0 \\
  \Rightarrow\; (\lambda_p - \lambda)(\lambda_r - \lambda) &= \mu\lambda_r\lambda ,
\end{align*}
exactly the PWE. However, we have disregarded the influence of the other waves.
Their contribution is given by
\begin{equation*}
  S_p := \sum_{n \neq p} \frac{\lambda}{\lambda_n^0 - \lambda} ,
\end{equation*}
leading to
\corr{Handwritten ending: ``$\ldots = 1 - \mu_\mathrm{eff} + \mu S_p$'', with
``$= \mu_\mathrm{eff}$'' over $\mu\,\lambda/(\lambda_p - \lambda)$. Removed:
$\mu\lambda/(\lambda_p - \lambda) \neq \mu_\mathrm{eff}$ (by the first relation
$\mu_\mathrm{eff} = 1/(\lambda/(\lambda_p-\lambda) + S_p)$), and the $S_p$
term has the wrong sign.}
\begin{equation}
  1 = \mu_\mathrm{eff}\left(\frac{\lambda}{\lambda_p - \lambda} + S_p\right)
  \;\Rightarrow\;
  1 - \beta = \mu\left(\frac{\lambda}{\lambda_p - \lambda} + S_p\right)
  \;\Rightarrow\;
  \beta = 1 - \mu\left(\frac{\lambda}{\lambda_p - \lambda} + S_p\right)
  \label{eq:beta_general}
\end{equation}
modifying the PWE.

As we are interested in an estimate for the upper edge of the gap, we consider
$\Gamma$, as in the effective medium picture ($S_p = 0$, only uniform wave)
band~2 starts at $f_r\sqrt{1+\mu}$ at $\mathbf{k} = 0$ and rises from there, so
$\Gamma$ is its minimum \fixed{in that picture. With $S$ included, $\Gamma$ is
a shallow local \emph{maximum} of band~2: moving off $\Gamma$, $S_p$ grows like
$|\mathbf{k}|^2$ while the uniform-wave term changes only like
$|\mathbf{k}|^4$, so the true minimum lies off $\Gamma$, $0.02$--$0.6\,\%$
lower, and the $\Gamma$ value is an upper bound for the edge}.%
\corr{Handwritten: ``so $\Gamma$ is its minimum''; true only for $S = 0$.}
Also, it makes our computation simpler. In the PWE, only terms close to or below
$\lambda$ have a significant contribution to the sum:
\begin{equation*}
  \frac{\lambda}{\lambda_n^0 - \lambda}
  \overset{x := \lambda_n^0/\lambda}{=}
  \frac{1}{x - 1}
  \;\to\;
  \left\{
  \begin{array}{l}
    \left.
    \begin{array}{ll}
      -1, & x \to 0 \\
      -\infty, & x \nearrow 1 \\
      +\infty, & x \searrow 1
    \end{array}
    \right\}\ \text{significant contribution} \\[2.2em]
    \begin{array}{ll}
      \;0, & x \to \infty
    \end{array}
  \end{array}
  \right.
\end{equation*}
For a plate whose\corr{Handwritten: ``who's''.} bare second band lies far above
the resonator's frequency at $\Gamma$, the only significant single contribution
is $\lambda_0^0 = 0$. We then have
\begin{equation*}
  \beta = 1 - \mu\left(\frac{\lambda}{0 - \lambda} + S_0\right)
  \equiv 1 - \mu(S - 1) = 1 + \mu(1 - S) \overset{!}{=} \frac{f^2}{f_r^2}
\end{equation*}
\begin{equation*}
  \Rightarrow\; f = f_r\sqrt{1 + \mu(1 - S)}
\end{equation*}
So the upper edge is estimated as the lower of two candidates:
\begin{itemize}
  \item The resonator band at about
        $f_\Gamma = f_r\sqrt{1 + \mu(1 - S)}$
  \item The sin-pattern not ``seeing'' the resonator at $X$ with frequency $f_X$.
\end{itemize}
\begin{equation}
  f_\text{upper} \approx \min\bigl(f_r\sqrt{1 + \mu(1 - S)},\; f_X\bigr)
  \label{eq:upper_edge_candidates}
\end{equation}
\corr{Equation~\eqref{eq:upper_edge_candidates} only restates the two bullet
points as a formula, so that the figures can refer to it.}

\begin{figure}[htbp]
  \centering
  \includegraphics[width=\textwidth]{Figures/upper_edge_candidates.pdf}
  \caption{The two upper-edge candidates of
    Eq.~\eqref{eq:upper_edge_candidates} against the tuning
    $\tau = f_r/f_X$, one panel per mass ratio (plane-wave model, centered
    resonator). Blue: the resonator band at $\Gamma$,
    $f_\Gamma = f_r\sqrt{1 + \mu(1 - S)}$, solid for the $0.1\,\mathrm{m}$
    cell ($h/a = 0.05$), dotted for Cool's cell ($h/a = 0.1$). Since $S$
    depends only on the plate and on $f$, this curve is exact and is drawn
    parametrically in $f$. Orange: the sin pattern at $X$, $f = f_X$. Dashed:
    the $S = 0$ estimate $f_r\sqrt{1+\mu}$ of
    Eq.~\eqref{eq:effective_medium_gap}; dotted grey: $f = f_r$. Black dots
    and open circles: the true upper and lower edges of the
    $0.1\,\mathrm{m}$ cell from the gap study of Fig.~\ref{fig:gap_map}. The
    candidates cross at
    $\tau_\text{switch} = 1/\sqrt{1 + \mu\bigl(1 - S(f_X)\bigr)}$
    (red dash-dotted line), with $S(f_X) = 0.424$ for the $0.1\,\mathrm{m}$
    cell and $0.514$ for Cool's cell. Below $\tau_\text{switch}$ the upper
    edge follows the resonator band, above it the edge is frozen at $f_X$.
    The true edge lies slightly below $f_\Gamma$ (by $0.02$--$0.6\,\%$),
    because $\Gamma$ is a local maximum of band~2 and its minimum lies off
    $\Gamma$. Last panel: $\tau_\text{switch}$ against $\mu$ for both cells
    and for $S = 0$.}
  \label{fig:upper_edge_candidates}
\end{figure}

\begin{figure}[htbp]
  \centering
  \includegraphics[width=\textwidth]{Figures/gap_map_cell0.1m_vs_cool0.05m.pdf}
  \caption{Complete band gaps against the tuning $\tau = f_r/f_X$ for five
    mass ratios (plane-wave model, centered resonator). Shaded: the
    $0.1\,\mathrm{m}$ cell ($h/a = 0.05$); dots: the $0.05\,\mathrm{m}$ cell
    with Cool's parameters ($h/a = 0.1$). Dashed: the effective-medium upper
    edge $f_r\sqrt{1+\mu}$ of Eq.~\eqref{eq:effective_medium_gap}; dotted:
    $f = f_r$. The upper edge follows the two candidates of
    Eq.~\eqref{eq:upper_edge_candidates}: for small $\tau$ it rises with the
    resonator band (below the dashed line, because of the bending sum $S$),
    for $\tau$ near and above $1$ it is frozen at $f_X$. Above $f_X$ the gap
    lies entirely below $f_r$ and closes once band~1 at $M$ reaches $f_X$
    (for $\mu = 1$ it does not close in the range shown). Red: the
    $1.2\,\mathrm{kHz}$ design (FE, $1031$--$1195\,\mathrm{Hz}$) of
    Fig.~\ref{fig:dispersion_with_il}.}
  \label{fig:gap_map}
\end{figure}

\begin{figure}[htbp]
  \centering
  \includegraphics[width=\textwidth]{Figures/gap_width_cell0.1m_vs_cool0.05m.pdf}
  \caption{Relative width of the widest complete gap against $\tau = f_r/f_X$
    (same models as Fig.~\ref{fig:gap_map}). Horizontal dashed lines: width of
    the effective-medium band $f_r \dots f_r\sqrt{1+\mu}$ of
    Eq.~\eqref{eq:effective_medium_gap}. The width is largest where the upper
    edge switches between the two candidates of
    Eq.~\eqref{eq:upper_edge_candidates}: below the switch both edges rise
    with $f_r$, above it the upper edge is frozen at $f_X$ while the lower
    edge keeps rising. The maxima lie at $\tau = 0.95, 0.95, 0.95, 0.90, 0.80$
    for $\mu = 0.05 \dots 1$, and the real gap there is about $1.5$ times wider
    than the effective-medium estimate.}
  \label{fig:gap_width}
\end{figure}

\begin{figure}[htbp]
  \centering
  \includegraphics[width=\textwidth]{Figures/insertion_loss_vs_gaps_1x1m_100resonators_mass0.25.pdf}
  \caption{Insertion loss of the radiated power (equal force) of the
    $1 \times 1\,\mathrm{m}$ plate with $10 \times 10$ resonators,
    $\mu = 0.25$, for three tunings, with the complete Bloch gap of the
    $0.1\,\mathrm{m}$ cell shaded. $f_r = 600\,\mathrm{Hz}$ ($\tau = 0.5$): the
    gap contains $f_r$ (classical local-resonance gap,
    Eq.~\eqref{eq:effective_medium_gap} with the bending-sum correction);
    $f_r = 1200\,\mathrm{Hz}$ ($\tau = 1$): the gap lies below $f_r$, upper
    edge $f_X$ (Fig.~\ref{fig:dispersion_with_il});
    $f_r = 2000\,\mathrm{Hz}$ ($\tau = 1.67$): no complete gap
    (Fig.~\ref{fig:gap_map}), and the IL stays weak. The IL peaks inside the
    gap in the first two cases and turns negative just above it, at the flat
    resonator band.}
  \label{fig:il_vs_gaps}
\end{figure}

\begin{figure}[htbp]
  \centering
  \includegraphics[width=\textwidth]{Figures/insertion_loss_shear_rule_all_tunings_1x1m_100resonators_mass0.25.pdf}
  \caption{The insertion loss of Fig.~\ref{fig:il_vs_gaps} computed with the
    reduced $2 \times 2$ and the full $3 \times 3$ shear integration. The
    reduced rule admits a spurious hourglass mode of the free plate, which
    distorts the response near the resonator band: with $2 \times 2$ the
    $1.2\,\mathrm{kHz}$ IL turns negative in the upper part of the gap, and the
    $600\,\mathrm{Hz}$ peak is lower. The forced-response pipeline uses
    $3 \times 3$.}
  \label{fig:il_shear_rule_all}
\end{figure}

\begin{figure}[htbp]
  \centering
  \includegraphics[width=\textwidth]{Figures/insertion_loss_shear_rule_1x1m_100resonators_1200hz_900to1500hz.pdf}
  \caption{Close-up of the $f_r = 1200\,\mathrm{Hz}$ case of
    Fig.~\ref{fig:il_shear_rule_all} ($900$--$1500\,\mathrm{Hz}$). With
    $3 \times 3$ the IL is positive across the whole Bloch gap (peak
    $+13.0\,\mathrm{dB}$ at $1120\,\mathrm{Hz}$) and reaches its minimum
    ($-13.9\,\mathrm{dB}$ at $1267\,\mathrm{Hz}$) at the flat resonator band,
    close to the resonator-band candidate $f_\Gamma$ of
    Eq.~\eqref{eq:upper_edge_candidates}.}
  \label{fig:il_shear_rule_1200}
\end{figure}

% =====================================================================

\providecommand{\fixed}[1]{{\color{blue!70!black}#1}}
\providecommand{\corr}[1]{\marginpar{\raggedright\footnotesize\color{red!70!black}#1}}

\chapter{Local resonator gap}
\label{ch:local_resonator_gap}

Let $m$ be the point mass, connected at $\mathbf{x}_0$ to the plate by a spring
of stiffness $k_s$. Write the plate deflection at the attachment point as
$w_0 := w(\mathbf{x}_0)$ and the displacement of the resonator mass as $u$.
The spring is stretched by $w_0 - u \Rightarrow m\ddot{u} = k_s(w_0 - u)$.
Assuming a harmonic motion $\ddot{u} = -\omega^2 u$, we obtain
\begin{equation*}
  -m\omega^2 u = k_s(w_0 - u)
  \;\Rightarrow\;
  u = \frac{k_s w_0}{k_s - m\omega^2}
  \;\overset{\cdot\frac{1/m}{1/m},\ \omega_r^2 = \frac{k_s}{m}}{=}\;
  \frac{\omega_r^2 w_0}{\omega_r^2 - \omega^2} .
\end{equation*}
The spring pulls on the plate with $F = k_s(u - w_0)$ according to Newton~III.
We find
\begin{equation*}
  F = k_s(u - w_0)
  \overset{k_s = m\omega_r^2}{=}
  m\omega_r^2\left(\frac{\omega_r^2}{\omega_r^2 - \omega^2} - 1\right) w_0
  = m\omega_r^2 \cdot \frac{\omega^2}{\omega_r^2 - \omega^2}\, w_0
  =: m_\mathrm{eff}(\omega)\,\omega^2 w_0 ,
\end{equation*}
where we have defined the \underline{effective mass}
$m_\mathrm{eff}(\omega) := m\,\dfrac{\omega_r^2}{\omega_r^2 - \omega^2}$.
Everything from Section~1, \emph{Bragg gaps}, now carries over by making the
replacements $m \to m_\mathrm{eff}(\omega)$ and
$\mu \to \mu_\mathrm{eff}(\omega) := \mu\,\dfrac{\omega_r^2}{\omega_r^2 - \omega^2}$.

We discuss different cases:
\begin{enumerate}
  \item $\omega \ll \omega_r$: $m_\mathrm{eff} \approx m$, $u \approx w_0$, like a
        rigid mass.
  \item $\omega$ just below $\omega_r$: $m_\mathrm{eff} \to \infty$,
        $u \gg w_0$, in phase, plate sees a huge mass and is pinned.
  \item $\omega$ just above $\omega_r$: $m_\mathrm{eff} \to -\infty$,
        $u \gg w_0$, antiphase, spring pulls plate back against its motion.
  \item $\omega \gg \omega_r$: $m_\mathrm{eff} \approx -m\,\dfrac{\omega_r^2}{\omega^2}$,
        $u \approx 0$, $F \to -k_s w_0 \mathrel{\widehat{=}}$ a spring to a fixed point.
\end{enumerate}

When considering \underline{waves much larger than the unit cell dimension $a$},
$m_\mathrm{eff}$ can be considered spread over the entire cell area:
$\rho h \to \rho h + m_\mathrm{eff}/A$.
Consider the Kirchhoff relation
\begin{equation*}
  D k^4 = \rho h\,(1 + \mu_\mathrm{eff})\,\omega^2
  = \rho h\left(1 + \mu\,\frac{\omega_r^2}{\omega_r^2 - \omega^2}\right)\omega^2
  = \rho h\,\frac{\omega_r^2(1+\mu) - \omega^2}{\omega_r^2 - \omega^2}\,\omega^2 .
\end{equation*}
A real solution needs a positive RHS, meaning that a negative RHS corresponds to
a band gap. As $\rho h, \omega^2 > 0$, consider the fraction: It is negative if
both the numerator $> 0$ and the denominator $< 0$ \fixed{(the opposite
combination, numerator $< 0$ and denominator $> 0$, would need
$\omega_r\sqrt{1+\mu} < \omega < \omega_r$, which is impossible)},%
\corr{Added: the second sign combination and why it is excluded.}
namely
\begin{equation*}
  \omega_r^2(1+\mu) - \omega^2 \overset{!}{>} 0
  \;\Rightarrow\; \omega < \omega_r\sqrt{1+\mu}
  \quad\text{and}\quad
  \omega_r \overset{!}{<} \omega ,
  \quad\text{thus}
\end{equation*}
\corr{Handwritten box: $\omega_r < \omega < \omega_r(1+\mu)$; the square root
was missing.}%
\begin{equation}
  \boxed{\;\omega_r < \omega < \fixed{\omega_r\sqrt{1+\mu}}\;}
  \label{eq:effective_medium_gap}
\end{equation}
This is limited by the fact that only $\mathbf{G} = 0$ is kept (i.e.,
$\tau = f_r/f_X \to 0$) and even at $\Gamma$ (small $\mathbf{k} \Rightarrow$
long waves) the upper edge is too high, as other waves add a local bending sum
$S$, lowering the upper edge by about $0.051\,\tau^2$\fixed{$\,f_X$ (fit for
$\mu = 0.25$)}.%
\corr{The $0.051\,\tau^2$ fit holds for $\mu = 0.25$ only and is in units
of $f_X$. Exact at $f_r = 1200\,\mathrm{Hz}$: $0.056\,f_X$
($1341.6 \to 1274.7\,\mathrm{Hz}$).}
Figure~\ref{fig:dispersion_with_il} shows how far the effective-medium band
\eqref{eq:effective_medium_gap} is off for the production design.

\begin{figure}[htbp]
  \centering
  \includegraphics[width=\textwidth]{Figures/dispersion_cell0.1m_1200hz_mass0.25_0to5000hz_with_il.pdf}
  \caption{Left: dispersion diagram (FE, Bloch) of the $0.1\,\mathrm{m}$ unit
    cell ($h = 5\,\mathrm{mm}$, steel) with one resonator,
    $f_r = 1200\,\mathrm{Hz}$, $\mu = 0.25$, along
    $\Gamma$--$X$--$M$--$\Gamma$; the color gives the plate share of the
    kinetic energy (blue: plate, red: resonator). The complete gap runs from
    $1031$ to $1195\,\mathrm{Hz}$ ($14.7\,\%$) and lies \emph{below} $f_r$,
    not in the effective-medium band
    $f_r \dots f_r\sqrt{1+\mu} = 1200\dots1342\,\mathrm{Hz}$ of
    Eq.~\eqref{eq:effective_medium_gap}: $f_r$ sits on the first Bragg fold
    ($f_X = 1194.4\,\mathrm{Hz}$). The two candidates for the upper edge of
    Eq.~\eqref{eq:upper_edge_candidates} are both visible: the flat resonator
    band at $\Gamma$ ($1273.5\,\mathrm{Hz}$, the effective-medium value lowered
    by the bending sum $S$) and the sin pattern at $X$ ($f_X$), which is the
    lower one here and sets the edge. The lower edge is band~1 at $M$. Right:
    insertion loss of the radiated power of the $1 \times 1\,\mathrm{m}$ plate
    with $10 \times 10$ resonators on the same $0.1\,\mathrm{m}$ lattice (same
    $f_r$ and $\mu$), for equal force and for equal input power. The IL peaks
    inside the gap and drops to its minimum just above it, at the flat
    resonator band.}
  \label{fig:dispersion_with_il}
\end{figure}

Keeping only one plate mode we obtain $w(\mathbf{x}) = q\,\varphi(\mathbf{x})$,
where $q$ is the modal amplitude. Let $u$ be the resonator's displacement and
$k_s$ its stiffness. Again assuming a harmonic response, i.e.,
$\ddot{u} = -\omega^2 u = -\lambda u$, we find
\begin{equation}
  k_s(u - \varphi q) = \lambda m u
  \label{eq:resonator_row}
\end{equation}
for the resonator. For the plate mode we find in analogy to the rigid mass
\begin{equation*}
  (\mathcal{L} - \lambda\rho h)\, w = f(\mathbf{x}) ,
  \qquad f(\mathbf{x}) = F\,\delta(\mathbf{x} - \mathbf{x}_0) .
\end{equation*}
Inserting $w = q\varphi$, multiplying by $\varphi^*$ and integrating over the
cell we again obtain (omission of the argument $\mathrel{\widehat{=}}$ evaluated
at $\mathbf{x}_0$)
\begin{equation}
  \lambda_p q - \lambda q = k_s\varphi^* u - k_s|\varphi|^2 q
  \;\Rightarrow\;
  (\lambda_p + k_s|\varphi|^2)\, q - k_s\varphi^* u = \lambda q .
  \label{eq:plate_row}
\end{equation}
Writing \eqref{eq:resonator_row} and \eqref{eq:plate_row} into matrix form
$K\mathbf{x} = \lambda M\mathbf{x}$, $\mathbf{x} := (q, u)^T$ we find
\begin{equation*}
  K = \begin{pmatrix}
        \lambda_p + k_s|\varphi|^2 & -k_s\varphi^* \\
        -k_s\varphi & k_s
      \end{pmatrix} ,
  \qquad
  M = \operatorname{diag}(1, m) .
\end{equation*}
Transforming the resonator coordinate into a mass-normalized coordinate, i.e.,
$v := \sqrt{m}\,u$ and identifying
\begin{gather*}
  k_s|\varphi|^2 = \frac{k_s}{\rho h A} = \frac{k_s}{m}\cdot\frac{m}{\rho h A}
  = \mu\lambda_r ; \qquad \frac{k_s}{m} = \lambda_r ; \\
  \frac{k_s\varphi}{\sqrt{m}} = \frac{k_s}{\sqrt{m\rho h A}}
  = \frac{k_s}{m}\sqrt{\frac{m}{\rho h A}} = \sqrt{\mu}\,\lambda_r
\end{gather*}
we see that $K$ becomes a Hermitian matrix $H$
\begin{equation*}
  H = \begin{pmatrix}
        \lambda_p + \mu\lambda_r & -\sqrt{\mu}\,\lambda_r \\
        -\sqrt{\mu}\,\lambda_r & \lambda_r
      \end{pmatrix}
\end{equation*}
and $M$ becomes the identity. The eigenvalues at equal diagonals, i.e.,
$\lambda_p = (1-\mu)\lambda_r$, are
\begin{equation*}
  \lambda_{1,2} = (1 \pm \sqrt{\mu})\,\lambda_r .
\end{equation*}

Now we want to consider the influence of other modes. For this, we rewrite
\eqref{eq:resonator_row} using
$\dfrac{m\omega^2}{k_s} = \dfrac{\omega^2}{\omega_r^2} =: \beta$ to find
\begin{equation*}
  u = \frac{w_0}{1 - \beta} .
\end{equation*}
Newton~III then gives us for the force on the
\fixed{plate}\corr{Handwritten: ``force on the mass''. $F = m\omega^2 u$ is
the force the mass exerts on the plate; the force on the mass is
$-m\omega^2 u$.}
\begin{equation*}
  F = m\omega^2 u = m\omega^2\,\frac{w_0}{1-\beta} = m_\mathrm{eff}\,\omega^2 w_0
  \qquad (\mu_\mathrm{eff}\text{ analogous}) .
\end{equation*}

As before, we find in our known notation
\begin{equation*}
  (\lambda_n^0 - \lambda)\, q_n = F\varphi_n^*(\mathbf{x}_0)
  \;\Rightarrow\;
  q_n = F\,\frac{\varphi_n^*(\mathbf{x}_0)}{\lambda_n^0 - \lambda} ,
\end{equation*}
i.e., how strongly each wave responds to the point force. We then find
\begin{equation*}
  w_0 = \sum_n \varphi_n(\mathbf{x}_0)\, q_n
  = F \sum_n \frac{|\varphi_n(\mathbf{x}_0)|^2}{\lambda_n^0 - \lambda}
  =: C F ,
\end{equation*}
where we have defined the compliance $C$, the
\fixed{deflection per unit force}.%
\corr{Handwritten: ``the force per unit deflection''. That is $Z$;
$C = w_0/F$ is the deflection per unit force.}
Defining $Z := m_\mathrm{eff}\lambda$, the force per unit deflection, we finally
obtain
\begin{equation*}
  w_0 = w_0 Z C \;\Rightarrow\; \boxed{1 = ZC} \qquad (w_0 \neq 0) .
\end{equation*}
Using $|\varphi(\mathbf{x}_0)|^2 = \dfrac{1}{\rho h A}$ we get
\begin{equation}
  \boxed{\;1 = \mu_\mathrm{eff}\sum_n \frac{\lambda}{\lambda_n^0 - \lambda}\;}
  \label{eq:pwe_condition}
\end{equation}
the plane-wave expansion (PWE) condition. At every $\mathbf{k}$, every
$\lambda$ satisfying it is one dot in the diagram. For $w_0 = 0$ we find
\begin{equation*}
  (\lambda_n^0 - \lambda)\, q_n = 0 \;\Rightarrow\; \lambda_n^0 = \lambda ,
\end{equation*}
thus, these modes keep their bare frequencies.

With this, we may find the upper edge of the band gap. First, writing out the
PWE for one \fixed{single}%
\corr{Handwritten: ``singular''.}
plate frequency $\lambda_p$, using
\begin{equation*}
  1 - \beta = 1 - \frac{\lambda}{\lambda_r} = \frac{\lambda_r - \lambda}{\lambda_r}
\end{equation*}
we get
\begin{equation*}
  1 = \mu\,\frac{\lambda_r}{\lambda_r - \lambda}\cdot\frac{\lambda}{\lambda_p - \lambda}
  \;\Rightarrow\;
  (\lambda_p - \lambda)(\lambda_r - \lambda) = \mu\lambda_r\lambda
\end{equation*}
and computing the characteristic polynomial of $H$ we find
\begin{align*}
  (\lambda_p + \mu\lambda_r - \lambda)(\lambda_r - \lambda) - \mu\lambda_r^2
  &= \bigl[(\lambda_p - \lambda) + \mu\lambda_r\bigr](\lambda_r - \lambda) - \mu\lambda_r^2 \\
  &= (\lambda_p - \lambda)(\lambda_r - \lambda) + \mu\lambda_r(\lambda_r - \lambda) - \mu\lambda_r^2
   = (\lambda_p - \lambda)(\lambda_r - \lambda) - \mu\lambda_r\lambda \overset{!}{=} 0 \\
  \Rightarrow\; (\lambda_p - \lambda)(\lambda_r - \lambda) &= \mu\lambda_r\lambda ,
\end{align*}
exactly the PWE. However, we have disregarded the influence of the other waves.
Their contribution is given by
\begin{equation*}
  S_p := \sum_{n \neq p} \frac{\lambda}{\lambda_n^0 - \lambda} ,
\end{equation*}
leading to
\corr{Handwritten ending: ``$\ldots = 1 - \mu_\mathrm{eff} + \mu S_p$'', with
``$= \mu_\mathrm{eff}$'' over $\mu\,\lambda/(\lambda_p - \lambda)$. Removed:
$\mu\lambda/(\lambda_p - \lambda) \neq \mu_\mathrm{eff}$ (by the first relation
$\mu_\mathrm{eff} = 1/(\lambda/(\lambda_p-\lambda) + S_p)$), and the $S_p$
term has the wrong sign.}
\begin{equation}
  1 = \mu_\mathrm{eff}\left(\frac{\lambda}{\lambda_p - \lambda} + S_p\right)
  \;\Rightarrow\;
  1 - \beta = \mu\left(\frac{\lambda}{\lambda_p - \lambda} + S_p\right)
  \;\Rightarrow\;
  \beta = 1 - \mu\left(\frac{\lambda}{\lambda_p - \lambda} + S_p\right)
  \label{eq:beta_general}
\end{equation}
modifying the PWE.

As we are interested in an estimate for the upper edge of the gap, we consider
$\Gamma$, as in the effective medium picture ($S_p = 0$, only uniform wave)
band~2 starts at $f_r\sqrt{1+\mu}$ at $\mathbf{k} = 0$ and rises from there, so
$\Gamma$ is its minimum \fixed{in that picture. With $S$ included, $\Gamma$ is
a shallow local \emph{maximum} of band~2: moving off $\Gamma$, $S_p$ grows like
$|\mathbf{k}|^2$ while the uniform-wave term changes only like
$|\mathbf{k}|^4$, so the true minimum lies off $\Gamma$, $0.02$--$0.6\,\%$
lower, and the $\Gamma$ value is an upper bound for the edge}.%
\corr{Handwritten: ``so $\Gamma$ is its minimum''; true only for $S = 0$.}
Also, it makes our computation simpler. In the PWE, only terms close to or below
$\lambda$ have a significant contribution to the sum:
\begin{equation*}
  \frac{\lambda}{\lambda_n^0 - \lambda}
  \overset{x := \lambda_n^0/\lambda}{=}
  \frac{1}{x - 1}
  \;\to\;
  \left\{
  \begin{array}{l}
    \left.
    \begin{array}{ll}
      -1, & x \to 0 \\
      -\infty, & x \nearrow 1 \\
      +\infty, & x \searrow 1
    \end{array}
    \right\}\ \text{significant contribution} \\[2.2em]
    \begin{array}{ll}
      \;0, & x \to \infty
    \end{array}
  \end{array}
  \right.
\end{equation*}
For a plate whose\corr{Handwritten: ``who's''.} bare second band lies far above
the resonator's frequency at $\Gamma$, the only significant single contribution
is $\lambda_0^0 = 0$. We then have
\begin{equation*}
  \beta = 1 - \mu\left(\frac{\lambda}{0 - \lambda} + S_0\right)
  \equiv 1 - \mu(S - 1) = 1 + \mu(1 - S) \overset{!}{=} \frac{f^2}{f_r^2}
\end{equation*}
\begin{equation*}
  \Rightarrow\; f = f_r\sqrt{1 + \mu(1 - S)}
\end{equation*}
So the upper edge is estimated as the lower of two candidates:
\begin{itemize}
  \item The resonator band at about
        $f_\Gamma = f_r\sqrt{1 + \mu(1 - S)}$
  \item The sin-pattern not ``seeing'' the resonator at $X$ with frequency $f_X$.
\end{itemize}
\begin{equation}
  f_\text{upper} \approx \min\bigl(f_r\sqrt{1 + \mu(1 - S)},\; f_X\bigr)
  \label{eq:upper_edge_candidates}
\end{equation}
\corr{Equation~\eqref{eq:upper_edge_candidates} only restates the two bullet
points as a formula, so that the figures can refer to it.}

\begin{figure}[htbp]
  \centering
  \includegraphics[width=\textwidth]{Figures/upper_edge_candidates.pdf}
  \caption{The two upper-edge candidates of
    Eq.~\eqref{eq:upper_edge_candidates} against the tuning
    $\tau = f_r/f_X$, one panel per mass ratio (plane-wave model, centered
    resonator). Blue: the resonator band at $\Gamma$,
    $f_\Gamma = f_r\sqrt{1 + \mu(1 - S)}$, solid for the $0.1\,\mathrm{m}$
    cell ($h/a = 0.05$), dotted for Cool's cell ($h/a = 0.1$). Since $S$
    depends only on the plate and on $f$, this curve is exact and is drawn
    parametrically in $f$. Orange: the sin pattern at $X$, $f = f_X$. Dashed:
    the $S = 0$ estimate $f_r\sqrt{1+\mu}$ of
    Eq.~\eqref{eq:effective_medium_gap}; dotted grey: $f = f_r$. Black dots
    and open circles: the true upper and lower edges of the
    $0.1\,\mathrm{m}$ cell from the gap study of Fig.~\ref{fig:gap_map}. The
    candidates cross at
    $\tau_\text{switch} = 1/\sqrt{1 + \mu\bigl(1 - S(f_X)\bigr)}$
    (red dash-dotted line), with $S(f_X) = 0.424$ for the $0.1\,\mathrm{m}$
    cell and $0.514$ for Cool's cell. Below $\tau_\text{switch}$ the upper
    edge follows the resonator band, above it the edge is frozen at $f_X$.
    The true edge lies slightly below $f_\Gamma$ (by $0.02$--$0.6\,\%$),
    because $\Gamma$ is a local maximum of band~2 and its minimum lies off
    $\Gamma$. Last panel: $\tau_\text{switch}$ against $\mu$ for both cells
    and for $S = 0$.}
  \label{fig:upper_edge_candidates}
\end{figure}

\begin{figure}[htbp]
  \centering
  \includegraphics[width=\textwidth]{Figures/gap_map_cell0.1m_vs_cool0.05m.pdf}
  \caption{Complete band gaps against the tuning $\tau = f_r/f_X$ for five
    mass ratios (plane-wave model, centered resonator). Shaded: the
    $0.1\,\mathrm{m}$ cell ($h/a = 0.05$); dots: the $0.05\,\mathrm{m}$ cell
    with Cool's parameters ($h/a = 0.1$). Dashed: the effective-medium upper
    edge $f_r\sqrt{1+\mu}$ of Eq.~\eqref{eq:effective_medium_gap}; dotted:
    $f = f_r$. The upper edge follows the two candidates of
    Eq.~\eqref{eq:upper_edge_candidates}: for small $\tau$ it rises with the
    resonator band (below the dashed line, because of the bending sum $S$),
    for $\tau$ near and above $1$ it is frozen at $f_X$. Above $f_X$ the gap
    lies entirely below $f_r$ and closes once band~1 at $M$ reaches $f_X$
    (for $\mu = 1$ it does not close in the range shown). Red: the
    $1.2\,\mathrm{kHz}$ design (FE, $1031$--$1195\,\mathrm{Hz}$) of
    Fig.~\ref{fig:dispersion_with_il}.}
  \label{fig:gap_map}
\end{figure}

\begin{figure}[htbp]
  \centering
  \includegraphics[width=\textwidth]{Figures/gap_width_cell0.1m_vs_cool0.05m.pdf}
  \caption{Relative width of the widest complete gap against $\tau = f_r/f_X$
    (same models as Fig.~\ref{fig:gap_map}). Horizontal dashed lines: width of
    the effective-medium band $f_r \dots f_r\sqrt{1+\mu}$ of
    Eq.~\eqref{eq:effective_medium_gap}. The width is largest where the upper
    edge switches between the two candidates of
    Eq.~\eqref{eq:upper_edge_candidates}: below the switch both edges rise
    with $f_r$, above it the upper edge is frozen at $f_X$ while the lower
    edge keeps rising. The maxima lie at $\tau = 0.95, 0.95, 0.95, 0.90, 0.80$
    for $\mu = 0.05 \dots 1$, and the real gap there is about $1.5$ times wider
    than the effective-medium estimate.}
  \label{fig:gap_width}
\end{figure}

\begin{figure}[htbp]
  \centering
  \includegraphics[width=\textwidth]{Figures/insertion_loss_vs_gaps_1x1m_100resonators_mass0.25.pdf}
  \caption{Insertion loss of the radiated power (equal force) of the
    $1 \times 1\,\mathrm{m}$ plate with $10 \times 10$ resonators,
    $\mu = 0.25$, for three tunings, with the complete Bloch gap of the
    $0.1\,\mathrm{m}$ cell shaded. $f_r = 600\,\mathrm{Hz}$ ($\tau = 0.5$): the
    gap contains $f_r$ (classical local-resonance gap,
    Eq.~\eqref{eq:effective_medium_gap} with the bending-sum correction);
    $f_r = 1200\,\mathrm{Hz}$ ($\tau = 1$): the gap lies below $f_r$, upper
    edge $f_X$ (Fig.~\ref{fig:dispersion_with_il});
    $f_r = 2000\,\mathrm{Hz}$ ($\tau = 1.67$): no complete gap
    (Fig.~\ref{fig:gap_map}), and the IL stays weak. The IL peaks inside the
    gap in the first two cases and turns negative just above it, at the flat
    resonator band.}
  \label{fig:il_vs_gaps}
\end{figure}

\begin{figure}[htbp]
  \centering
  \includegraphics[width=\textwidth]{Figures/insertion_loss_shear_rule_all_tunings_1x1m_100resonators_mass0.25.pdf}
  \caption{The insertion loss of Fig.~\ref{fig:il_vs_gaps} computed with the
    reduced $2 \times 2$ and the full $3 \times 3$ shear integration. The
    reduced rule admits a spurious hourglass mode of the free plate, which
    distorts the response near the resonator band: with $2 \times 2$ the
    $1.2\,\mathrm{kHz}$ IL turns negative in the upper part of the gap, and the
    $600\,\mathrm{Hz}$ peak is lower. The forced-response pipeline uses
    $3 \times 3$.}
  \label{fig:il_shear_rule_all}
\end{figure}

\begin{figure}[htbp]
  \centering
  \includegraphics[width=\textwidth]{Figures/insertion_loss_shear_rule_1x1m_100resonators_1200hz_900to1500hz.pdf}
  \caption{Close-up of the $f_r = 1200\,\mathrm{Hz}$ case of
    Fig.~\ref{fig:il_shear_rule_all} ($900$--$1500\,\mathrm{Hz}$). With
    $3 \times 3$ the IL is positive across the whole Bloch gap (peak
    $+13.0\,\mathrm{dB}$ at $1120\,\mathrm{Hz}$) and reaches its minimum
    ($-13.9\,\mathrm{dB}$ at $1267\,\mathrm{Hz}$) at the flat resonator band,
    close to the resonator-band candidate $f_\Gamma$ of
    Eq.~\eqref{eq:upper_edge_candidates}.}
  \label{fig:il_shear_rule_1200}
\end{figure}

% =====================================================================

\providecommand{\fixed}[1]{{\color{blue!70!black}#1}}
\providecommand{\corr}[1]{\marginpar{\raggedright\footnotesize\color{red!70!black}#1}}

\chapter{Local resonator gap}
\label{ch:local_resonator_gap}

Let $m$ be the point mass, connected at $\mathbf{x}_0$ to the plate by a spring
of stiffness $k_s$. Write the plate deflection at the attachment point as
$w_0 := w(\mathbf{x}_0)$ and the displacement of the resonator mass as $u$.
The spring is stretched by $w_0 - u \Rightarrow m\ddot{u} = k_s(w_0 - u)$.
Assuming a harmonic motion $\ddot{u} = -\omega^2 u$, we obtain
\begin{equation*}
  -m\omega^2 u = k_s(w_0 - u)
  \;\Rightarrow\;
  u = \frac{k_s w_0}{k_s - m\omega^2}
  \;\overset{\cdot\frac{1/m}{1/m},\ \omega_r^2 = \frac{k_s}{m}}{=}\;
  \frac{\omega_r^2 w_0}{\omega_r^2 - \omega^2} .
\end{equation*}
The spring pulls on the plate with $F = k_s(u - w_0)$ according to Newton~III.
We find
\begin{equation*}
  F = k_s(u - w_0)
  \overset{k_s = m\omega_r^2}{=}
  m\omega_r^2\left(\frac{\omega_r^2}{\omega_r^2 - \omega^2} - 1\right) w_0
  = m\omega_r^2 \cdot \frac{\omega^2}{\omega_r^2 - \omega^2}\, w_0
  =: m_\mathrm{eff}(\omega)\,\omega^2 w_0 ,
\end{equation*}
where we have defined the \underline{effective mass}
$m_\mathrm{eff}(\omega) := m\,\dfrac{\omega_r^2}{\omega_r^2 - \omega^2}$.
Everything from Section~1, \emph{Bragg gaps}, now carries over by making the
replacements $m \to m_\mathrm{eff}(\omega)$ and
$\mu \to \mu_\mathrm{eff}(\omega) := \mu\,\dfrac{\omega_r^2}{\omega_r^2 - \omega^2}$.

We discuss different cases:
\begin{enumerate}
  \item $\omega \ll \omega_r$: $m_\mathrm{eff} \approx m$, $u \approx w_0$, like a
        rigid mass.
  \item $\omega$ just below $\omega_r$: $m_\mathrm{eff} \to \infty$,
        $u \gg w_0$, in phase, plate sees a huge mass and is pinned.
  \item $\omega$ just above $\omega_r$: $m_\mathrm{eff} \to -\infty$,
        $u \gg w_0$, antiphase, spring pulls plate back against its motion.
  \item $\omega \gg \omega_r$: $m_\mathrm{eff} \approx -m\,\dfrac{\omega_r^2}{\omega^2}$,
        $u \approx 0$, $F \to -k_s w_0 \mathrel{\widehat{=}}$ a spring to a fixed point.
\end{enumerate}

When considering \underline{waves much larger than the unit cell dimension $a$},
$m_\mathrm{eff}$ can be considered spread over the entire cell area:
$\rho h \to \rho h + m_\mathrm{eff}/A$.
Consider the Kirchhoff relation
\begin{equation*}
  D k^4 = \rho h\,(1 + \mu_\mathrm{eff})\,\omega^2
  = \rho h\left(1 + \mu\,\frac{\omega_r^2}{\omega_r^2 - \omega^2}\right)\omega^2
  = \rho h\,\frac{\omega_r^2(1+\mu) - \omega^2}{\omega_r^2 - \omega^2}\,\omega^2 .
\end{equation*}
A real solution needs a positive RHS, meaning that a negative RHS corresponds to
a band gap. As $\rho h, \omega^2 > 0$, consider the fraction: It is negative if
both the numerator $> 0$ and the denominator $< 0$ \fixed{(the opposite
combination, numerator $< 0$ and denominator $> 0$, would need
$\omega_r\sqrt{1+\mu} < \omega < \omega_r$, which is impossible)},%
\corr{Added: the second sign combination and why it is excluded.}
namely
\begin{equation*}
  \omega_r^2(1+\mu) - \omega^2 \overset{!}{>} 0
  \;\Rightarrow\; \omega < \omega_r\sqrt{1+\mu}
  \quad\text{and}\quad
  \omega_r \overset{!}{<} \omega ,
  \quad\text{thus}
\end{equation*}
\corr{Handwritten box: $\omega_r < \omega < \omega_r(1+\mu)$; the square root
was missing.}%
\begin{equation}
  \boxed{\;\omega_r < \omega < \fixed{\omega_r\sqrt{1+\mu}}\;}
  \label{eq:effective_medium_gap}
\end{equation}
This is limited by the fact that only $\mathbf{G} = 0$ is kept (i.e.,
$\tau = f_r/f_X \to 0$) and even at $\Gamma$ (small $\mathbf{k} \Rightarrow$
long waves) the upper edge is too high, as other waves add a local bending sum
$S$, lowering the upper edge by about $0.051\,\tau^2$\fixed{$\,f_X$ (fit for
$\mu = 0.25$)}.%
\corr{The $0.051\,\tau^2$ fit holds for $\mu = 0.25$ only and is in units
of $f_X$. Exact at $f_r = 1200\,\mathrm{Hz}$: $0.056\,f_X$
($1341.6 \to 1274.7\,\mathrm{Hz}$).}
Figure~\ref{fig:dispersion_with_il} shows how far the effective-medium band
\eqref{eq:effective_medium_gap} is off for the production design.

\begin{figure}[htbp]
  \centering
  \includegraphics[width=\textwidth]{Figures/dispersion_cell0.1m_1200hz_mass0.25_0to5000hz_with_il.pdf}
  \caption{Left: dispersion diagram (FE, Bloch) of the $0.1\,\mathrm{m}$ unit
    cell ($h = 5\,\mathrm{mm}$, steel) with one resonator,
    $f_r = 1200\,\mathrm{Hz}$, $\mu = 0.25$, along
    $\Gamma$--$X$--$M$--$\Gamma$; the color gives the plate share of the
    kinetic energy (blue: plate, red: resonator). The complete gap runs from
    $1031$ to $1195\,\mathrm{Hz}$ ($14.7\,\%$) and lies \emph{below} $f_r$,
    not in the effective-medium band
    $f_r \dots f_r\sqrt{1+\mu} = 1200\dots1342\,\mathrm{Hz}$ of
    Eq.~\eqref{eq:effective_medium_gap}: $f_r$ sits on the first Bragg fold
    ($f_X = 1194.4\,\mathrm{Hz}$). The two candidates for the upper edge of
    Eq.~\eqref{eq:upper_edge_candidates} are both visible: the flat resonator
    band at $\Gamma$ ($1273.5\,\mathrm{Hz}$, the effective-medium value lowered
    by the bending sum $S$) and the sin pattern at $X$ ($f_X$), which is the
    lower one here and sets the edge. The lower edge is band~1 at $M$. Right:
    insertion loss of the radiated power of the $1 \times 1\,\mathrm{m}$ plate
    with $10 \times 10$ resonators on the same $0.1\,\mathrm{m}$ lattice (same
    $f_r$ and $\mu$), for equal force and for equal input power. The IL peaks
    inside the gap and drops to its minimum just above it, at the flat
    resonator band.}
  \label{fig:dispersion_with_il}
\end{figure}

Keeping only one plate mode we obtain $w(\mathbf{x}) = q\,\varphi(\mathbf{x})$,
where $q$ is the modal amplitude. Let $u$ be the resonator's displacement and
$k_s$ its stiffness. Again assuming a harmonic response, i.e.,
$\ddot{u} = -\omega^2 u = -\lambda u$, we find
\begin{equation}
  k_s(u - \varphi q) = \lambda m u
  \label{eq:resonator_row}
\end{equation}
for the resonator. For the plate mode we find in analogy to the rigid mass
\begin{equation*}
  (\mathcal{L} - \lambda\rho h)\, w = f(\mathbf{x}) ,
  \qquad f(\mathbf{x}) = F\,\delta(\mathbf{x} - \mathbf{x}_0) .
\end{equation*}
Inserting $w = q\varphi$, multiplying by $\varphi^*$ and integrating over the
cell we again obtain (omission of the argument $\mathrel{\widehat{=}}$ evaluated
at $\mathbf{x}_0$)
\begin{equation}
  \lambda_p q - \lambda q = k_s\varphi^* u - k_s|\varphi|^2 q
  \;\Rightarrow\;
  (\lambda_p + k_s|\varphi|^2)\, q - k_s\varphi^* u = \lambda q .
  \label{eq:plate_row}
\end{equation}
Writing \eqref{eq:resonator_row} and \eqref{eq:plate_row} into matrix form
$K\mathbf{x} = \lambda M\mathbf{x}$, $\mathbf{x} := (q, u)^T$ we find
\begin{equation*}
  K = \begin{pmatrix}
        \lambda_p + k_s|\varphi|^2 & -k_s\varphi^* \\
        -k_s\varphi & k_s
      \end{pmatrix} ,
  \qquad
  M = \operatorname{diag}(1, m) .
\end{equation*}
Transforming the resonator coordinate into a mass-normalized coordinate, i.e.,
$v := \sqrt{m}\,u$ and identifying
\begin{gather*}
  k_s|\varphi|^2 = \frac{k_s}{\rho h A} = \frac{k_s}{m}\cdot\frac{m}{\rho h A}
  = \mu\lambda_r ; \qquad \frac{k_s}{m} = \lambda_r ; \\
  \frac{k_s\varphi}{\sqrt{m}} = \frac{k_s}{\sqrt{m\rho h A}}
  = \frac{k_s}{m}\sqrt{\frac{m}{\rho h A}} = \sqrt{\mu}\,\lambda_r
\end{gather*}
we see that $K$ becomes a Hermitian matrix $H$
\begin{equation*}
  H = \begin{pmatrix}
        \lambda_p + \mu\lambda_r & -\sqrt{\mu}\,\lambda_r \\
        -\sqrt{\mu}\,\lambda_r & \lambda_r
      \end{pmatrix}
\end{equation*}
and $M$ becomes the identity. The eigenvalues at equal diagonals, i.e.,
$\lambda_p = (1-\mu)\lambda_r$, are
\begin{equation*}
  \lambda_{1,2} = (1 \pm \sqrt{\mu})\,\lambda_r .
\end{equation*}

Now we want to consider the influence of other modes. For this, we rewrite
\eqref{eq:resonator_row} using
$\dfrac{m\omega^2}{k_s} = \dfrac{\omega^2}{\omega_r^2} =: \beta$ to find
\begin{equation*}
  u = \frac{w_0}{1 - \beta} .
\end{equation*}
Newton~III then gives us for the force on the
\fixed{plate}\corr{Handwritten: ``force on the mass''. $F = m\omega^2 u$ is
the force the mass exerts on the plate; the force on the mass is
$-m\omega^2 u$.}
\begin{equation*}
  F = m\omega^2 u = m\omega^2\,\frac{w_0}{1-\beta} = m_\mathrm{eff}\,\omega^2 w_0
  \qquad (\mu_\mathrm{eff}\text{ analogous}) .
\end{equation*}

As before, we find in our known notation
\begin{equation*}
  (\lambda_n^0 - \lambda)\, q_n = F\varphi_n^*(\mathbf{x}_0)
  \;\Rightarrow\;
  q_n = F\,\frac{\varphi_n^*(\mathbf{x}_0)}{\lambda_n^0 - \lambda} ,
\end{equation*}
i.e., how strongly each wave responds to the point force. We then find
\begin{equation*}
  w_0 = \sum_n \varphi_n(\mathbf{x}_0)\, q_n
  = F \sum_n \frac{|\varphi_n(\mathbf{x}_0)|^2}{\lambda_n^0 - \lambda}
  =: C F ,
\end{equation*}
where we have defined the compliance $C$, the
\fixed{deflection per unit force}.%
\corr{Handwritten: ``the force per unit deflection''. That is $Z$;
$C = w_0/F$ is the deflection per unit force.}
Defining $Z := m_\mathrm{eff}\lambda$, the force per unit deflection, we finally
obtain
\begin{equation*}
  w_0 = w_0 Z C \;\Rightarrow\; \boxed{1 = ZC} \qquad (w_0 \neq 0) .
\end{equation*}
Using $|\varphi(\mathbf{x}_0)|^2 = \dfrac{1}{\rho h A}$ we get
\begin{equation}
  \boxed{\;1 = \mu_\mathrm{eff}\sum_n \frac{\lambda}{\lambda_n^0 - \lambda}\;}
  \label{eq:pwe_condition}
\end{equation}
the plane-wave expansion (PWE) condition. At every $\mathbf{k}$, every
$\lambda$ satisfying it is one dot in the diagram. For $w_0 = 0$ we find
\begin{equation*}
  (\lambda_n^0 - \lambda)\, q_n = 0 \;\Rightarrow\; \lambda_n^0 = \lambda ,
\end{equation*}
thus, these modes keep their bare frequencies.

With this, we may find the upper edge of the band gap. First, writing out the
PWE for one \fixed{single}%
\corr{Handwritten: ``singular''.}
plate frequency $\lambda_p$, using
\begin{equation*}
  1 - \beta = 1 - \frac{\lambda}{\lambda_r} = \frac{\lambda_r - \lambda}{\lambda_r}
\end{equation*}
we get
\begin{equation*}
  1 = \mu\,\frac{\lambda_r}{\lambda_r - \lambda}\cdot\frac{\lambda}{\lambda_p - \lambda}
  \;\Rightarrow\;
  (\lambda_p - \lambda)(\lambda_r - \lambda) = \mu\lambda_r\lambda
\end{equation*}
and computing the characteristic polynomial of $H$ we find
\begin{align*}
  (\lambda_p + \mu\lambda_r - \lambda)(\lambda_r - \lambda) - \mu\lambda_r^2
  &= \bigl[(\lambda_p - \lambda) + \mu\lambda_r\bigr](\lambda_r - \lambda) - \mu\lambda_r^2 \\
  &= (\lambda_p - \lambda)(\lambda_r - \lambda) + \mu\lambda_r(\lambda_r - \lambda) - \mu\lambda_r^2
   = (\lambda_p - \lambda)(\lambda_r - \lambda) - \mu\lambda_r\lambda \overset{!}{=} 0 \\
  \Rightarrow\; (\lambda_p - \lambda)(\lambda_r - \lambda) &= \mu\lambda_r\lambda ,
\end{align*}
exactly the PWE. However, we have disregarded the influence of the other waves.
Their contribution is given by
\begin{equation*}
  S_p := \sum_{n \neq p} \frac{\lambda}{\lambda_n^0 - \lambda} ,
\end{equation*}
leading to
\corr{Handwritten ending: ``$\ldots = 1 - \mu_\mathrm{eff} + \mu S_p$'', with
``$= \mu_\mathrm{eff}$'' over $\mu\,\lambda/(\lambda_p - \lambda)$. Removed:
$\mu\lambda/(\lambda_p - \lambda) \neq \mu_\mathrm{eff}$ (by the first relation
$\mu_\mathrm{eff} = 1/(\lambda/(\lambda_p-\lambda) + S_p)$), and the $S_p$
term has the wrong sign.}
\begin{equation}
  1 = \mu_\mathrm{eff}\left(\frac{\lambda}{\lambda_p - \lambda} + S_p\right)
  \;\Rightarrow\;
  1 - \beta = \mu\left(\frac{\lambda}{\lambda_p - \lambda} + S_p\right)
  \;\Rightarrow\;
  \beta = 1 - \mu\left(\frac{\lambda}{\lambda_p - \lambda} + S_p\right)
  \label{eq:beta_general}
\end{equation}
modifying the PWE.

As we are interested in an estimate for the upper edge of the gap, we consider
$\Gamma$, as in the effective medium picture ($S_p = 0$, only uniform wave)
band~2 starts at $f_r\sqrt{1+\mu}$ at $\mathbf{k} = 0$ and rises from there, so
$\Gamma$ is its minimum \fixed{in that picture. With $S$ included, $\Gamma$ is
a shallow local \emph{maximum} of band~2: moving off $\Gamma$, $S_p$ grows like
$|\mathbf{k}|^2$ while the uniform-wave term changes only like
$|\mathbf{k}|^4$, so the true minimum lies off $\Gamma$, $0.02$--$0.6\,\%$
lower, and the $\Gamma$ value is an upper bound for the edge}.%
\corr{Handwritten: ``so $\Gamma$ is its minimum''; true only for $S = 0$.}
Also, it makes our computation simpler. In the PWE, only terms close to or below
$\lambda$ have a significant contribution to the sum:
\begin{equation*}
  \frac{\lambda}{\lambda_n^0 - \lambda}
  \overset{x := \lambda_n^0/\lambda}{=}
  \frac{1}{x - 1}
  \;\to\;
  \left\{
  \begin{array}{l}
    \left.
    \begin{array}{ll}
      -1, & x \to 0 \\
      -\infty, & x \nearrow 1 \\
      +\infty, & x \searrow 1
    \end{array}
    \right\}\ \text{significant contribution} \\[2.2em]
    \begin{array}{ll}
      \;0, & x \to \infty
    \end{array}
  \end{array}
  \right.
\end{equation*}
For a plate whose\corr{Handwritten: ``who's''.} bare second band lies far above
the resonator's frequency at $\Gamma$, the only significant single contribution
is $\lambda_0^0 = 0$. We then have
\begin{equation*}
  \beta = 1 - \mu\left(\frac{\lambda}{0 - \lambda} + S_0\right)
  \equiv 1 - \mu(S - 1) = 1 + \mu(1 - S) \overset{!}{=} \frac{f^2}{f_r^2}
\end{equation*}
\begin{equation*}
  \Rightarrow\; f = f_r\sqrt{1 + \mu(1 - S)}
\end{equation*}
So the upper edge is estimated as the lower of two candidates:
\begin{itemize}
  \item The resonator band at about
        $f_\Gamma = f_r\sqrt{1 + \mu(1 - S)}$
  \item The sin-pattern not ``seeing'' the resonator at $X$ with frequency $f_X$.
\end{itemize}
\begin{equation}
  f_\text{upper} \approx \min\bigl(f_r\sqrt{1 + \mu(1 - S)},\; f_X\bigr)
  \label{eq:upper_edge_candidates}
\end{equation}
\corr{Equation~\eqref{eq:upper_edge_candidates} only restates the two bullet
points as a formula, so that the figures can refer to it.}

\begin{figure}[htbp]
  \centering
  \includegraphics[width=\textwidth]{Figures/upper_edge_candidates.pdf}
  \caption{The two upper-edge candidates of
    Eq.~\eqref{eq:upper_edge_candidates} against the tuning
    $\tau = f_r/f_X$, one panel per mass ratio (plane-wave model, centered
    resonator). Blue: the resonator band at $\Gamma$,
    $f_\Gamma = f_r\sqrt{1 + \mu(1 - S)}$, solid for the $0.1\,\mathrm{m}$
    cell ($h/a = 0.05$), dotted for Cool's cell ($h/a = 0.1$). Since $S$
    depends only on the plate and on $f$, this curve is exact and is drawn
    parametrically in $f$. Orange: the sin pattern at $X$, $f = f_X$. Dashed:
    the $S = 0$ estimate $f_r\sqrt{1+\mu}$ of
    Eq.~\eqref{eq:effective_medium_gap}; dotted grey: $f = f_r$. Black dots
    and open circles: the true upper and lower edges of the
    $0.1\,\mathrm{m}$ cell from the gap study of Fig.~\ref{fig:gap_map}. The
    candidates cross at
    $\tau_\text{switch} = 1/\sqrt{1 + \mu\bigl(1 - S(f_X)\bigr)}$
    (red dash-dotted line), with $S(f_X) = 0.424$ for the $0.1\,\mathrm{m}$
    cell and $0.514$ for Cool's cell. Below $\tau_\text{switch}$ the upper
    edge follows the resonator band, above it the edge is frozen at $f_X$.
    The true edge lies slightly below $f_\Gamma$ (by $0.02$--$0.6\,\%$),
    because $\Gamma$ is a local maximum of band~2 and its minimum lies off
    $\Gamma$. Last panel: $\tau_\text{switch}$ against $\mu$ for both cells
    and for $S = 0$.}
  \label{fig:upper_edge_candidates}
\end{figure}

\begin{figure}[htbp]
  \centering
  \includegraphics[width=\textwidth]{Figures/gap_map_cell0.1m_vs_cool0.05m.pdf}
  \caption{Complete band gaps against the tuning $\tau = f_r/f_X$ for five
    mass ratios (plane-wave model, centered resonator). Shaded: the
    $0.1\,\mathrm{m}$ cell ($h/a = 0.05$); dots: the $0.05\,\mathrm{m}$ cell
    with Cool's parameters ($h/a = 0.1$). Dashed: the effective-medium upper
    edge $f_r\sqrt{1+\mu}$ of Eq.~\eqref{eq:effective_medium_gap}; dotted:
    $f = f_r$. The upper edge follows the two candidates of
    Eq.~\eqref{eq:upper_edge_candidates}: for small $\tau$ it rises with the
    resonator band (below the dashed line, because of the bending sum $S$),
    for $\tau$ near and above $1$ it is frozen at $f_X$. Above $f_X$ the gap
    lies entirely below $f_r$ and closes once band~1 at $M$ reaches $f_X$
    (for $\mu = 1$ it does not close in the range shown). Red: the
    $1.2\,\mathrm{kHz}$ design (FE, $1031$--$1195\,\mathrm{Hz}$) of
    Fig.~\ref{fig:dispersion_with_il}.}
  \label{fig:gap_map}
\end{figure}

\begin{figure}[htbp]
  \centering
  \includegraphics[width=\textwidth]{Figures/gap_width_cell0.1m_vs_cool0.05m.pdf}
  \caption{Relative width of the widest complete gap against $\tau = f_r/f_X$
    (same models as Fig.~\ref{fig:gap_map}). Horizontal dashed lines: width of
    the effective-medium band $f_r \dots f_r\sqrt{1+\mu}$ of
    Eq.~\eqref{eq:effective_medium_gap}. The width is largest where the upper
    edge switches between the two candidates of
    Eq.~\eqref{eq:upper_edge_candidates}: below the switch both edges rise
    with $f_r$, above it the upper edge is frozen at $f_X$ while the lower
    edge keeps rising. The maxima lie at $\tau = 0.95, 0.95, 0.95, 0.90, 0.80$
    for $\mu = 0.05 \dots 1$, and the real gap there is about $1.5$ times wider
    than the effective-medium estimate.}
  \label{fig:gap_width}
\end{figure}

\begin{figure}[htbp]
  \centering
  \includegraphics[width=\textwidth]{Figures/insertion_loss_vs_gaps_1x1m_100resonators_mass0.25.pdf}
  \caption{Insertion loss of the radiated power (equal force) of the
    $1 \times 1\,\mathrm{m}$ plate with $10 \times 10$ resonators,
    $\mu = 0.25$, for three tunings, with the complete Bloch gap of the
    $0.1\,\mathrm{m}$ cell shaded. $f_r = 600\,\mathrm{Hz}$ ($\tau = 0.5$): the
    gap contains $f_r$ (classical local-resonance gap,
    Eq.~\eqref{eq:effective_medium_gap} with the bending-sum correction);
    $f_r = 1200\,\mathrm{Hz}$ ($\tau = 1$): the gap lies below $f_r$, upper
    edge $f_X$ (Fig.~\ref{fig:dispersion_with_il});
    $f_r = 2000\,\mathrm{Hz}$ ($\tau = 1.67$): no complete gap
    (Fig.~\ref{fig:gap_map}), and the IL stays weak. The IL peaks inside the
    gap in the first two cases and turns negative just above it, at the flat
    resonator band.}
  \label{fig:il_vs_gaps}
\end{figure}

\begin{figure}[htbp]
  \centering
  \includegraphics[width=\textwidth]{Figures/insertion_loss_shear_rule_all_tunings_1x1m_100resonators_mass0.25.pdf}
  \caption{The insertion loss of Fig.~\ref{fig:il_vs_gaps} computed with the
    reduced $2 \times 2$ and the full $3 \times 3$ shear integration. The
    reduced rule admits a spurious hourglass mode of the free plate, which
    distorts the response near the resonator band: with $2 \times 2$ the
    $1.2\,\mathrm{kHz}$ IL turns negative in the upper part of the gap, and the
    $600\,\mathrm{Hz}$ peak is lower. The forced-response pipeline uses
    $3 \times 3$.}
  \label{fig:il_shear_rule_all}
\end{figure}

\begin{figure}[htbp]
  \centering
  \includegraphics[width=\textwidth]{Figures/insertion_loss_shear_rule_1x1m_100resonators_1200hz_900to1500hz.pdf}
  \caption{Close-up of the $f_r = 1200\,\mathrm{Hz}$ case of
    Fig.~\ref{fig:il_shear_rule_all} ($900$--$1500\,\mathrm{Hz}$). With
    $3 \times 3$ the IL is positive across the whole Bloch gap (peak
    $+13.0\,\mathrm{dB}$ at $1120\,\mathrm{Hz}$) and reaches its minimum
    ($-13.9\,\mathrm{dB}$ at $1267\,\mathrm{Hz}$) at the flat resonator band,
    close to the resonator-band candidate $f_\Gamma$ of
    Eq.~\eqref{eq:upper_edge_candidates}.}
  \label{fig:il_shear_rule_1200}
\end{figure}

% =====================================================================

\providecommand{\fixed}[1]{{\color{blue!70!black}#1}}
\providecommand{\corr}[1]{\marginpar{\raggedright\footnotesize\color{red!70!black}#1}}

\chapter{Local resonator gap}
\label{ch:local_resonator_gap}

Let $m$ be the point mass, connected at $\mathbf{x}_0$ to the plate by a spring
of stiffness $k_s$. Write the plate deflection at the attachment point as
$w_0 := w(\mathbf{x}_0)$ and the displacement of the resonator mass as $u$.
The spring is stretched by $w_0 - u \Rightarrow m\ddot{u} = k_s(w_0 - u)$.
Assuming a harmonic motion $\ddot{u} = -\omega^2 u$, we obtain
\begin{equation*}
  -m\omega^2 u = k_s(w_0 - u)
  \;\Rightarrow\;
  u = \frac{k_s w_0}{k_s - m\omega^2}
  \;\overset{\cdot\frac{1/m}{1/m},\ \omega_r^2 = \frac{k_s}{m}}{=}\;
  \frac{\omega_r^2 w_0}{\omega_r^2 - \omega^2} .
\end{equation*}
The spring pulls on the plate with $F = k_s(u - w_0)$ according to Newton~III.
We find
\begin{equation*}
  F = k_s(u - w_0)
  \overset{k_s = m\omega_r^2}{=}
  m\omega_r^2\left(\frac{\omega_r^2}{\omega_r^2 - \omega^2} - 1\right) w_0
  = m\omega_r^2 \cdot \frac{\omega^2}{\omega_r^2 - \omega^2}\, w_0
  =: m_\mathrm{eff}(\omega)\,\omega^2 w_0 ,
\end{equation*}
where we have defined the \underline{effective mass}
$m_\mathrm{eff}(\omega) := m\,\dfrac{\omega_r^2}{\omega_r^2 - \omega^2}$.
Everything from Section~1, \emph{Bragg gaps}, now carries over by making the
replacements $m \to m_\mathrm{eff}(\omega)$ and
$\mu \to \mu_\mathrm{eff}(\omega) := \mu\,\dfrac{\omega_r^2}{\omega_r^2 - \omega^2}$.

We discuss different cases:
\begin{enumerate}
  \item $\omega \ll \omega_r$: $m_\mathrm{eff} \approx m$, $u \approx w_0$, like a
        rigid mass.
  \item $\omega$ just below $\omega_r$: $m_\mathrm{eff} \to \infty$,
        $u \gg w_0$, in phase, plate sees a huge mass and is pinned.
  \item $\omega$ just above $\omega_r$: $m_\mathrm{eff} \to -\infty$,
        $u \gg w_0$, antiphase, spring pulls plate back against its motion.
  \item $\omega \gg \omega_r$: $m_\mathrm{eff} \approx -m\,\dfrac{\omega_r^2}{\omega^2}$,
        $u \approx 0$, $F \to -k_s w_0 \mathrel{\widehat{=}}$ a spring to a fixed point.
\end{enumerate}

When considering \underline{waves much larger than the unit cell dimension $a$},
$m_\mathrm{eff}$ can be considered spread over the entire cell area:
$\rho h \to \rho h + m_\mathrm{eff}/A$.
Consider the Kirchhoff relation
\begin{equation*}
  D k^4 = \rho h\,(1 + \mu_\mathrm{eff})\,\omega^2
  = \rho h\left(1 + \mu\,\frac{\omega_r^2}{\omega_r^2 - \omega^2}\right)\omega^2
  = \rho h\,\frac{\omega_r^2(1+\mu) - \omega^2}{\omega_r^2 - \omega^2}\,\omega^2 .
\end{equation*}
A real solution needs a positive RHS, meaning that a negative RHS corresponds to
a band gap. As $\rho h, \omega^2 > 0$, consider the fraction: It is negative if
both the numerator $> 0$ and the denominator $< 0$ \fixed{(the opposite
combination, numerator $< 0$ and denominator $> 0$, would need
$\omega_r\sqrt{1+\mu} < \omega < \omega_r$, which is impossible)},%
\corr{Added: the second sign combination and why it is excluded.}
namely
\begin{equation*}
  \omega_r^2(1+\mu) - \omega^2 \overset{!}{>} 0
  \;\Rightarrow\; \omega < \omega_r\sqrt{1+\mu}
  \quad\text{and}\quad
  \omega_r \overset{!}{<} \omega ,
  \quad\text{thus}
\end{equation*}
\corr{Handwritten box: $\omega_r < \omega < \omega_r(1+\mu)$; the square root
was missing.}%
\begin{equation}
  \boxed{\;\omega_r < \omega < \fixed{\omega_r\sqrt{1+\mu}}\;}
  \label{eq:effective_medium_gap}
\end{equation}
This is limited by the fact that only $\mathbf{G} = 0$ is kept (i.e.,
$\tau = f_r/f_X \to 0$) and even at $\Gamma$ (small $\mathbf{k} \Rightarrow$
long waves) the upper edge is too high, as other waves add a local bending sum
$S$, lowering the upper edge by about $0.051\,\tau^2$\fixed{$\,f_X$ (fit for
$\mu = 0.25$)}.%
\corr{The $0.051\,\tau^2$ fit holds for $\mu = 0.25$ only and is in units
of $f_X$. Exact at $f_r = 1200\,\mathrm{Hz}$: $0.056\,f_X$
($1341.6 \to 1274.7\,\mathrm{Hz}$).}
Figure~\ref{fig:dispersion_with_il} shows how far the effective-medium band
\eqref{eq:effective_medium_gap} is off for the production design.

\begin{figure}[htbp]
  \centering
  \includegraphics[width=\textwidth]{Figures/dispersion_cell0.1m_1200hz_mass0.25_0to5000hz_with_il.pdf}
  \caption{Left: dispersion diagram (FE, Bloch) of the $0.1\,\mathrm{m}$ unit
    cell ($h = 5\,\mathrm{mm}$, steel) with one resonator,
    $f_r = 1200\,\mathrm{Hz}$, $\mu = 0.25$, along
    $\Gamma$--$X$--$M$--$\Gamma$; the color gives the plate share of the
    kinetic energy (blue: plate, red: resonator). The complete gap runs from
    $1031$ to $1195\,\mathrm{Hz}$ ($14.7\,\%$) and lies \emph{below} $f_r$,
    not in the effective-medium band
    $f_r \dots f_r\sqrt{1+\mu} = 1200\dots1342\,\mathrm{Hz}$ of
    Eq.~\eqref{eq:effective_medium_gap}: $f_r$ sits on the first Bragg fold
    ($f_X = 1194.4\,\mathrm{Hz}$). The two candidates for the upper edge of
    Eq.~\eqref{eq:upper_edge_candidates} are both visible: the flat resonator
    band at $\Gamma$ ($1273.5\,\mathrm{Hz}$, the effective-medium value lowered
    by the bending sum $S$) and the sin pattern at $X$ ($f_X$), which is the
    lower one here and sets the edge. The lower edge is band~1 at $M$. Right:
    insertion loss of the radiated power of the $1 \times 1\,\mathrm{m}$ plate
    with $10 \times 10$ resonators on the same $0.1\,\mathrm{m}$ lattice (same
    $f_r$ and $\mu$), for equal force and for equal input power. The IL peaks
    inside the gap and drops to its minimum just above it, at the flat
    resonator band.}
  \label{fig:dispersion_with_il}
\end{figure}

Keeping only one plate mode we obtain $w(\mathbf{x}) = q\,\varphi(\mathbf{x})$,
where $q$ is the modal amplitude. Let $u$ be the resonator's displacement and
$k_s$ its stiffness. Again assuming a harmonic response, i.e.,
$\ddot{u} = -\omega^2 u = -\lambda u$, we find
\begin{equation}
  k_s(u - \varphi q) = \lambda m u
  \label{eq:resonator_row}
\end{equation}
for the resonator. For the plate mode we find in analogy to the rigid mass
\begin{equation*}
  (\mathcal{L} - \lambda\rho h)\, w = f(\mathbf{x}) ,
  \qquad f(\mathbf{x}) = F\,\delta(\mathbf{x} - \mathbf{x}_0) .
\end{equation*}
Inserting $w = q\varphi$, multiplying by $\varphi^*$ and integrating over the
cell we again obtain (omission of the argument $\mathrel{\widehat{=}}$ evaluated
at $\mathbf{x}_0$)
\begin{equation}
  \lambda_p q - \lambda q = k_s\varphi^* u - k_s|\varphi|^2 q
  \;\Rightarrow\;
  (\lambda_p + k_s|\varphi|^2)\, q - k_s\varphi^* u = \lambda q .
  \label{eq:plate_row}
\end{equation}
Writing \eqref{eq:resonator_row} and \eqref{eq:plate_row} into matrix form
$K\mathbf{x} = \lambda M\mathbf{x}$, $\mathbf{x} := (q, u)^T$ we find
\begin{equation*}
  K = \begin{pmatrix}
        \lambda_p + k_s|\varphi|^2 & -k_s\varphi^* \\
        -k_s\varphi & k_s
      \end{pmatrix} ,
  \qquad
  M = \operatorname{diag}(1, m) .
\end{equation*}
Transforming the resonator coordinate into a mass-normalized coordinate, i.e.,
$v := \sqrt{m}\,u$ and identifying
\begin{gather*}
  k_s|\varphi|^2 = \frac{k_s}{\rho h A} = \frac{k_s}{m}\cdot\frac{m}{\rho h A}
  = \mu\lambda_r ; \qquad \frac{k_s}{m} = \lambda_r ; \\
  \frac{k_s\varphi}{\sqrt{m}} = \frac{k_s}{\sqrt{m\rho h A}}
  = \frac{k_s}{m}\sqrt{\frac{m}{\rho h A}} = \sqrt{\mu}\,\lambda_r
\end{gather*}
we see that $K$ becomes a Hermitian matrix $H$
\begin{equation*}
  H = \begin{pmatrix}
        \lambda_p + \mu\lambda_r & -\sqrt{\mu}\,\lambda_r \\
        -\sqrt{\mu}\,\lambda_r & \lambda_r
      \end{pmatrix}
\end{equation*}
and $M$ becomes the identity. The eigenvalues at equal diagonals, i.e.,
$\lambda_p = (1-\mu)\lambda_r$, are
\begin{equation*}
  \lambda_{1,2} = (1 \pm \sqrt{\mu})\,\lambda_r .
\end{equation*}

Now we want to consider the influence of other modes. For this, we rewrite
\eqref{eq:resonator_row} using
$\dfrac{m\omega^2}{k_s} = \dfrac{\omega^2}{\omega_r^2} =: \beta$ to find
\begin{equation*}
  u = \frac{w_0}{1 - \beta} .
\end{equation*}
Newton~III then gives us for the force on the
\fixed{plate}\corr{Handwritten: ``force on the mass''. $F = m\omega^2 u$ is
the force the mass exerts on the plate; the force on the mass is
$-m\omega^2 u$.}
\begin{equation*}
  F = m\omega^2 u = m\omega^2\,\frac{w_0}{1-\beta} = m_\mathrm{eff}\,\omega^2 w_0
  \qquad (\mu_\mathrm{eff}\text{ analogous}) .
\end{equation*}

As before, we find in our known notation
\begin{equation*}
  (\lambda_n^0 - \lambda)\, q_n = F\varphi_n^*(\mathbf{x}_0)
  \;\Rightarrow\;
  q_n = F\,\frac{\varphi_n^*(\mathbf{x}_0)}{\lambda_n^0 - \lambda} ,
\end{equation*}
i.e., how strongly each wave responds to the point force. We then find
\begin{equation*}
  w_0 = \sum_n \varphi_n(\mathbf{x}_0)\, q_n
  = F \sum_n \frac{|\varphi_n(\mathbf{x}_0)|^2}{\lambda_n^0 - \lambda}
  =: C F ,
\end{equation*}
where we have defined the compliance $C$, the
\fixed{deflection per unit force}.%
\corr{Handwritten: ``the force per unit deflection''. That is $Z$;
$C = w_0/F$ is the deflection per unit force.}
Defining $Z := m_\mathrm{eff}\lambda$, the force per unit deflection, we finally
obtain
\begin{equation*}
  w_0 = w_0 Z C \;\Rightarrow\; \boxed{1 = ZC} \qquad (w_0 \neq 0) .
\end{equation*}
Using $|\varphi(\mathbf{x}_0)|^2 = \dfrac{1}{\rho h A}$ we get
\begin{equation}
  \boxed{\;1 = \mu_\mathrm{eff}\sum_n \frac{\lambda}{\lambda_n^0 - \lambda}\;}
  \label{eq:pwe_condition}
\end{equation}
the plane-wave expansion (PWE) condition. At every $\mathbf{k}$, every
$\lambda$ satisfying it is one dot in the diagram. For $w_0 = 0$ we find
\begin{equation*}
  (\lambda_n^0 - \lambda)\, q_n = 0 \;\Rightarrow\; \lambda_n^0 = \lambda ,
\end{equation*}
thus, these modes keep their bare frequencies.

With this, we may find the upper edge of the band gap. First, writing out the
PWE for one \fixed{single}%
\corr{Handwritten: ``singular''.}
plate frequency $\lambda_p$, using
\begin{equation*}
  1 - \beta = 1 - \frac{\lambda}{\lambda_r} = \frac{\lambda_r - \lambda}{\lambda_r}
\end{equation*}
we get
\begin{equation*}
  1 = \mu\,\frac{\lambda_r}{\lambda_r - \lambda}\cdot\frac{\lambda}{\lambda_p - \lambda}
  \;\Rightarrow\;
  (\lambda_p - \lambda)(\lambda_r - \lambda) = \mu\lambda_r\lambda
\end{equation*}
and computing the characteristic polynomial of $H$ we find
\begin{align*}
  (\lambda_p + \mu\lambda_r - \lambda)(\lambda_r - \lambda) - \mu\lambda_r^2
  &= \bigl[(\lambda_p - \lambda) + \mu\lambda_r\bigr](\lambda_r - \lambda) - \mu\lambda_r^2 \\
  &= (\lambda_p - \lambda)(\lambda_r - \lambda) + \mu\lambda_r(\lambda_r - \lambda) - \mu\lambda_r^2
   = (\lambda_p - \lambda)(\lambda_r - \lambda) - \mu\lambda_r\lambda \overset{!}{=} 0 \\
  \Rightarrow\; (\lambda_p - \lambda)(\lambda_r - \lambda) &= \mu\lambda_r\lambda ,
\end{align*}
exactly the PWE. However, we have disregarded the influence of the other waves.
Their contribution is given by
\begin{equation*}
  S_p := \sum_{n \neq p} \frac{\lambda}{\lambda_n^0 - \lambda} ,
\end{equation*}
leading to
\corr{Handwritten ending: ``$\ldots = 1 - \mu_\mathrm{eff} + \mu S_p$'', with
``$= \mu_\mathrm{eff}$'' over $\mu\,\lambda/(\lambda_p - \lambda)$. Removed:
$\mu\lambda/(\lambda_p - \lambda) \neq \mu_\mathrm{eff}$ (by the first relation
$\mu_\mathrm{eff} = 1/(\lambda/(\lambda_p-\lambda) + S_p)$), and the $S_p$
term has the wrong sign.}
\begin{equation}
  1 = \mu_\mathrm{eff}\left(\frac{\lambda}{\lambda_p - \lambda} + S_p\right)
  \;\Rightarrow\;
  1 - \beta = \mu\left(\frac{\lambda}{\lambda_p - \lambda} + S_p\right)
  \;\Rightarrow\;
  \beta = 1 - \mu\left(\frac{\lambda}{\lambda_p - \lambda} + S_p\right)
  \label{eq:beta_general}
\end{equation}
modifying the PWE.

As we are interested in an estimate for the upper edge of the gap, we consider
$\Gamma$, as in the effective medium picture ($S_p = 0$, only uniform wave)
band~2 starts at $f_r\sqrt{1+\mu}$ at $\mathbf{k} = 0$ and rises from there, so
$\Gamma$ is its minimum \fixed{in that picture. With $S$ included, $\Gamma$ is
a shallow local \emph{maximum} of band~2: moving off $\Gamma$, $S_p$ grows like
$|\mathbf{k}|^2$ while the uniform-wave term changes only like
$|\mathbf{k}|^4$, so the true minimum lies off $\Gamma$, $0.02$--$0.6\,\%$
lower, and the $\Gamma$ value is an upper bound for the edge}.%
\corr{Handwritten: ``so $\Gamma$ is its minimum''; true only for $S = 0$.}
Also, it makes our computation simpler. In the PWE, only terms close to or below
$\lambda$ have a significant contribution to the sum:
\begin{equation*}
  \frac{\lambda}{\lambda_n^0 - \lambda}
  \overset{x := \lambda_n^0/\lambda}{=}
  \frac{1}{x - 1}
  \;\to\;
  \left\{
  \begin{array}{l}
    \left.
    \begin{array}{ll}
      -1, & x \to 0 \\
      -\infty, & x \nearrow 1 \\
      +\infty, & x \searrow 1
    \end{array}
    \right\}\ \text{significant contribution} \\[2.2em]
    \begin{array}{ll}
      \;0, & x \to \infty
    \end{array}
  \end{array}
  \right.
\end{equation*}
For a plate whose\corr{Handwritten: ``who's''.} bare second band lies far above
the resonator's frequency at $\Gamma$, the only significant single contribution
is $\lambda_0^0 = 0$. We then have
\begin{equation*}
  \beta = 1 - \mu\left(\frac{\lambda}{0 - \lambda} + S_0\right)
  \equiv 1 - \mu(S - 1) = 1 + \mu(1 - S) \overset{!}{=} \frac{f^2}{f_r^2}
\end{equation*}
\begin{equation*}
  \Rightarrow\; f = f_r\sqrt{1 + \mu(1 - S)}
\end{equation*}
So the upper edge is estimated as the lower of two candidates:
\begin{itemize}
  \item The resonator band at about
        $f_\Gamma = f_r\sqrt{1 + \mu(1 - S)}$
  \item The sin-pattern not ``seeing'' the resonator at $X$ with frequency $f_X$.
\end{itemize}
\begin{equation}
  f_\text{upper} \approx \min\bigl(f_r\sqrt{1 + \mu(1 - S)},\; f_X\bigr)
  \label{eq:upper_edge_candidates}
\end{equation}
\corr{Equation~\eqref{eq:upper_edge_candidates} only restates the two bullet
points as a formula, so that the figures can refer to it.}

\begin{figure}[htbp]
  \centering
  \includegraphics[width=\textwidth]{Figures/upper_edge_candidates.pdf}
  \caption{The two upper-edge candidates of
    Eq.~\eqref{eq:upper_edge_candidates} against the tuning
    $\tau = f_r/f_X$, one panel per mass ratio (plane-wave model, centered
    resonator). Blue: the resonator band at $\Gamma$,
    $f_\Gamma = f_r\sqrt{1 + \mu(1 - S)}$, solid for the $0.1\,\mathrm{m}$
    cell ($h/a = 0.05$), dotted for Cool's cell ($h/a = 0.1$). Since $S$
    depends only on the plate and on $f$, this curve is exact and is drawn
    parametrically in $f$. Orange: the sin pattern at $X$, $f = f_X$. Dashed:
    the $S = 0$ estimate $f_r\sqrt{1+\mu}$ of
    Eq.~\eqref{eq:effective_medium_gap}; dotted grey: $f = f_r$. Black dots
    and open circles: the true upper and lower edges of the
    $0.1\,\mathrm{m}$ cell from the gap study of Fig.~\ref{fig:gap_map}. The
    candidates cross at
    $\tau_\text{switch} = 1/\sqrt{1 + \mu\bigl(1 - S(f_X)\bigr)}$
    (red dash-dotted line), with $S(f_X) = 0.424$ for the $0.1\,\mathrm{m}$
    cell and $0.514$ for Cool's cell. Below $\tau_\text{switch}$ the upper
    edge follows the resonator band, above it the edge is frozen at $f_X$.
    The true edge lies slightly below $f_\Gamma$ (by $0.02$--$0.6\,\%$),
    because $\Gamma$ is a local maximum of band~2 and its minimum lies off
    $\Gamma$. Last panel: $\tau_\text{switch}$ against $\mu$ for both cells
    and for $S = 0$.}
  \label{fig:upper_edge_candidates}
\end{figure}

\begin{figure}[htbp]
  \centering
  \includegraphics[width=\textwidth]{Figures/gap_map_cell0.1m_vs_cool0.05m.pdf}
  \caption{Complete band gaps against the tuning $\tau = f_r/f_X$ for five
    mass ratios (plane-wave model, centered resonator). Shaded: the
    $0.1\,\mathrm{m}$ cell ($h/a = 0.05$); dots: the $0.05\,\mathrm{m}$ cell
    with Cool's parameters ($h/a = 0.1$). Dashed: the effective-medium upper
    edge $f_r\sqrt{1+\mu}$ of Eq.~\eqref{eq:effective_medium_gap}; dotted:
    $f = f_r$. The upper edge follows the two candidates of
    Eq.~\eqref{eq:upper_edge_candidates}: for small $\tau$ it rises with the
    resonator band (below the dashed line, because of the bending sum $S$),
    for $\tau$ near and above $1$ it is frozen at $f_X$. Above $f_X$ the gap
    lies entirely below $f_r$ and closes once band~1 at $M$ reaches $f_X$
    (for $\mu = 1$ it does not close in the range shown). Red: the
    $1.2\,\mathrm{kHz}$ design (FE, $1031$--$1195\,\mathrm{Hz}$) of
    Fig.~\ref{fig:dispersion_with_il}.}
  \label{fig:gap_map}
\end{figure}

\begin{figure}[htbp]
  \centering
  \includegraphics[width=\textwidth]{Figures/gap_width_cell0.1m_vs_cool0.05m.pdf}
  \caption{Relative width of the widest complete gap against $\tau = f_r/f_X$
    (same models as Fig.~\ref{fig:gap_map}). Horizontal dashed lines: width of
    the effective-medium band $f_r \dots f_r\sqrt{1+\mu}$ of
    Eq.~\eqref{eq:effective_medium_gap}. The width is largest where the upper
    edge switches between the two candidates of
    Eq.~\eqref{eq:upper_edge_candidates}: below the switch both edges rise
    with $f_r$, above it the upper edge is frozen at $f_X$ while the lower
    edge keeps rising. The maxima lie at $\tau = 0.95, 0.95, 0.95, 0.90, 0.80$
    for $\mu = 0.05 \dots 1$, and the real gap there is about $1.5$ times wider
    than the effective-medium estimate.}
  \label{fig:gap_width}
\end{figure}

\begin{figure}[htbp]
  \centering
  \includegraphics[width=\textwidth]{Figures/insertion_loss_vs_gaps_1x1m_100resonators_mass0.25.pdf}
  \caption{Insertion loss of the radiated power (equal force) of the
    $1 \times 1\,\mathrm{m}$ plate with $10 \times 10$ resonators,
    $\mu = 0.25$, for three tunings, with the complete Bloch gap of the
    $0.1\,\mathrm{m}$ cell shaded. $f_r = 600\,\mathrm{Hz}$ ($\tau = 0.5$): the
    gap contains $f_r$ (classical local-resonance gap,
    Eq.~\eqref{eq:effective_medium_gap} with the bending-sum correction);
    $f_r = 1200\,\mathrm{Hz}$ ($\tau = 1$): the gap lies below $f_r$, upper
    edge $f_X$ (Fig.~\ref{fig:dispersion_with_il});
    $f_r = 2000\,\mathrm{Hz}$ ($\tau = 1.67$): no complete gap
    (Fig.~\ref{fig:gap_map}), and the IL stays weak. The IL peaks inside the
    gap in the first two cases and turns negative just above it, at the flat
    resonator band.}
  \label{fig:il_vs_gaps}
\end{figure}

\begin{figure}[htbp]
  \centering
  \includegraphics[width=\textwidth]{Figures/insertion_loss_shear_rule_all_tunings_1x1m_100resonators_mass0.25.pdf}
  \caption{The insertion loss of Fig.~\ref{fig:il_vs_gaps} computed with the
    reduced $2 \times 2$ and the full $3 \times 3$ shear integration. The
    reduced rule admits a spurious hourglass mode of the free plate, which
    distorts the response near the resonator band: with $2 \times 2$ the
    $1.2\,\mathrm{kHz}$ IL turns negative in the upper part of the gap, and the
    $600\,\mathrm{Hz}$ peak is lower. The forced-response pipeline uses
    $3 \times 3$.}
  \label{fig:il_shear_rule_all}
\end{figure}

\begin{figure}[htbp]
  \centering
  \includegraphics[width=\textwidth]{Figures/insertion_loss_shear_rule_1x1m_100resonators_1200hz_900to1500hz.pdf}
  \caption{Close-up of the $f_r = 1200\,\mathrm{Hz}$ case of
    Fig.~\ref{fig:il_shear_rule_all} ($900$--$1500\,\mathrm{Hz}$). With
    $3 \times 3$ the IL is positive across the whole Bloch gap (peak
    $+13.0\,\mathrm{dB}$ at $1120\,\mathrm{Hz}$) and reaches its minimum
    ($-13.9\,\mathrm{dB}$ at $1267\,\mathrm{Hz}$) at the flat resonator band,
    close to the resonator-band candidate $f_\Gamma$ of
    Eq.~\eqref{eq:upper_edge_candidates}.}
  \label{fig:il_shear_rule_1200}
\end{figure}
